\documentclass[12pt,a4paper]{ctexbook}
\usepackage{geometry}
\geometry{margin=2.2cm}
\usepackage{setspace}
\setstretch{1.35}
\usepackage{hyperref}
\title{唐诗三百首}
\author{古典文献整理}
\date{}
\begin{document}
\maketitle
\tableofcontents
\newpage
\section{第1则}
行宮·元稹\par
寥落古行宮，宮花寂寞紅。\par
白頭宮女在，閒坐說玄宗。\par
登鸛雀樓·王之渙\par
白日依山盡，黃河入海流。\par
欲窮千里目，更上一層樓。\par
新嫁娘詞·王建\par
三日入廚下，洗手作羹湯。\par
未諳姑食性，先遣小姑嘗。\par
相思·王維\par
紅豆生南國，春來發幾枝。\par
願君多采擷，此物最相思。\par
雜詩·王維\par
君自故鄉來，應知故鄉事。\par
來日綺窗前，寒梅著花未？\par
鹿柴·王維\par
空山不見人，但聞人語響。\par
返景入深林，復照青苔上。\par
竹裏館·王維\par
獨坐幽篁裏，彈琴復長嘯。\par
深林人不知，明月來相照。\par
山中送別·王維\par
山中相送罷，日暮掩柴扉。\par
春草明年綠，王孫歸不歸？(明年 一作：年年)\par
問劉十九·白居易\par
綠蟻新醅酒，紅泥小火爐。\par
晚來天欲雪，能飲一杯無？\par
哥舒歌·西鄙人\par
北斗七星高，哥舒夜帶刀。\par
至今窺牧馬，不敢過臨洮。\par
靜夜思·李白\par
床前明月光，疑是地上霜，\par
舉頭望明月，低頭思故鄉。\par
怨情·李白\par
美人卷珠簾，深坐顰蛾眉。(顰 一作：蹙)\par
但見淚痕溼，不知心恨誰。\par
登樂遊原·李商隱\par
向晚意不適，驅車登古原。\par
夕陽無限好，只是近黃昏。\par
聽箏·李端\par
鳴箏金粟柱，素手玉房前。\par
欲得周郎顧，時時誤拂弦。\par
渡漢江·宋之問\par
嶺外音書斷，經冬復歷春。\par
近鄉情更怯，不敢問來人。\par
八陣圖·杜甫\par
功蓋三分國，名成八陣圖。（名成 一作：名高）\par
江流石不轉，遺恨失吞吳。\par
宿建德江·孟浩然\par
移舟泊煙渚，日暮客愁新。\par
野曠天低樹，江清月近人。\par
春曉·孟浩然\par
春眠不覺曉，處處聞啼鳥。\par
夜來風雨聲，花落知多少。\par
春怨·金昌緒\par
打起黃鶯兒，莫教枝上啼。\par
啼時驚妾夢，不得到遼西。\par
江雪·柳宗元\par
千山鳥飛絕，萬徑人蹤滅。\par
孤舟蓑笠翁，獨釣寒江雪。\par
秋夜寄邱員外·韋應物\par
懷君屬秋夜，散步詠涼天。\par
空山松子落，幽人應未眠。\par
終南望餘雪·祖詠\par
終南陰嶺秀，積雪浮雲端。\par
林表明霽色，城中增暮寒。\par
宮詞·張祜\par
故國三千里，深宮二十年。\par
一聲何滿子，雙淚落君前。\par
尋隱者不遇·賈島\par
松下問童子，言師採藥去。\par
只在此山中，雲深不知處。\par
送崔九·裴迪\par
歸山深淺去，須盡丘壑美。\par
莫學武陵人，暫遊桃源裏。\par
送靈澈上人·劉長卿\par
蒼蒼竹林寺，杳杳鐘聲晚。\par
荷笠帶斜陽，青山獨歸遠。(斜陽 一作：夕陽)\par
聽彈琴·劉長卿\par
泠泠七絃上，靜聽松風寒。(七絃 一作：七絲)\par
古調雖自愛，今人多不彈。\par
送上人·劉長卿\par
孤雲將野鶴，豈向人間住。\par
莫買沃洲山，時人已知處。\par
玉臺體·權德輿\par
昨夜裙帶解，今朝蟢子飛。\par
鉛華不可棄，莫是藁砧歸。\par

\section{第2则}
芙蓉樓送辛漸·王昌齡\par
寒雨連江夜入吳，平明送客楚山孤。\par
洛陽親友如相問，一片冰心在玉壺。\par
閨怨·王昌齡\par
閨中少婦不知愁，春日凝妝上翠樓。\par
忽見陌頭楊柳色，悔教夫婿覓封侯。\par
春宮曲·王昌齡\par
昨夜風開露井桃，未央前殿月輪高。\par
平陽歌舞新承寵，簾外春寒賜錦袍。\par
九月九日憶山東兄弟·王維\par
獨在異鄉爲異客，每逢佳節倍思親。\par
遙知兄弟登高處，遍插茱萸少一人。\par
涼州詞·王翰\par
葡萄美酒夜光杯，欲飲琵琶馬上催。\par
醉臥沙場君莫笑，古來征戰幾人回？\par
後宮詞·白居易\par
淚溼羅巾夢不成，夜深前殿按歌聲。\par
紅顏未老恩先斷，斜倚薰籠坐到明。\par
宮中詞·朱慶餘\par
寂寂花時閉院門，美人相並立瓊軒。\par
含情欲說宮中事，鸚鵡前頭不敢言。\par
近試上張水部·朱慶餘\par
洞房昨夜停紅燭，待曉堂前拜舅姑。\par
妝罷低聲問夫婿，畫眉深淺入時無。\par
逢入京使·岑參\par
故園東望路漫漫，雙袖龍鍾淚不幹。\par
馬上相逢無紙筆，憑君傳語報平安。\par
黃鶴樓送孟浩然之廣陵·李白\par
故人西辭黃鶴樓，煙花三月下揚州。\par
孤帆遠影碧空盡，唯見長江天際流。 (唯 通：惟)\par
早發白帝城·李白\par
朝辭白帝彩雲間，千里江陵一日還。\par
兩岸猿聲啼不住，輕舟已過萬重山。\par
夜上受降城聞笛·李益\par
回樂烽前沙似雪，受降城外月如霜。(回樂\par
一作：回樂\par
）\par
不知何處吹蘆管，一夜徵人盡望鄉。\par
賈生·李商隱\par
宣室求賢訪逐臣，賈生才調更無倫。\par
可憐夜半虛前席，不問蒼生問鬼神。\par
隋宮·李商隱\par
乘興南遊不戒嚴，九重誰省諫書函。\par
春風舉國裁宮錦，半作障泥半作帆。\par
瑤池·李商隱\par
瑤池阿母綺窗開，黃竹歌聲動地哀。\par
八駿日行三萬裏，穆王何事不重來。\par
嫦娥·李商隱\par
雲母屏風燭影深，長河漸落曉星沉。\par
嫦娥應悔偷靈藥，碧海青天夜夜心。\par
夜雨寄北·李商隱\par
君問歸期未有期，巴山夜雨漲秋池。\par
何當共剪西窗燭，卻話巴山夜雨時。\par
寄令狐郎中·李商隱\par
嵩雲秦樹久離居，雙鯉迢迢一紙書。\par
休問梁園舊賓客，茂陵秋雨病相如。\par
爲有·李商隱\par
爲有云屏無限嬌，鳳城寒盡怕春宵。\par
無端嫁得金龜婿，辜負香衾事早朝。\par
江南逢李龜年·杜甫\par
岐王宅裏尋常見，崔九堂前幾度聞。\par
正是江南好風景，落花時節又逢君。\par
贈別·其一·杜牧\par
娉娉嫋嫋十三餘，豆蔻梢頭二月初。\par
春風十裏揚州路，卷上珠簾總不如。\par
贈別·其二·杜牧\par
多情卻似總無情，唯覺樽前笑不成。\par
蠟燭有心還惜別，替人垂淚到天明。\par
金谷園·杜牧\par
繁華事散逐香塵，流水無情草自春。\par
日暮東風怨啼鳥，落花猶似墜樓人。\par
寄揚州韓綽判官·杜牧\par
青山隱隱水迢迢，秋盡江南草未凋。\par
二十四橋明月夜，玉人何處教吹簫？\par
遣懷·杜牧\par
落魄江南載酒行，楚腰纖細掌中輕。(江南 一作：江湖；纖細 一作：腸斷)\par
十年一覺揚州夢，贏得青樓薄倖名。\par
秋夕·杜牧\par
銀燭秋光冷畫屏，輕羅小扇撲流螢。\par
天階夜色涼如水，臥看牽牛織女星。\par
將赴吳興登樂遊原一絕·杜牧\par
清時有味是無能，閒愛孤雲靜愛僧。\par
欲把一麾江海去，樂遊原上望昭陵。\par
赤壁·杜牧\par
折戟沉沙鐵未銷，自將磨洗認前朝。\par
東風不與周郎便，銅雀春深鎖二喬。\par
泊秦淮·杜牧\par
煙籠寒水月籠沙，夜泊秦淮近酒家。\par
商女不知亡國恨，隔江猶唱後庭花。\par
徵人怨·柳中庸\par
歲歲金河復玉關，朝朝馬策與刀環。\par
三春白雪歸青冢，萬里黃河繞黑山。\par
金陵圖·韋莊\par
誰謂傷心畫不成，畫人心逐世人情。\par
君看六幅南朝事，老木寒雲滿故城。\par
滁州西澗·韋應物\par
獨憐幽草澗邊生，上有黃鸝深樹鳴。\par
春潮帶雨晚來急，野渡無人舟自橫。\par
桃花溪·張旭\par
隱隱飛橋隔野煙，石磯西畔問漁船。\par
桃花盡日隨流水，洞在清溪何處邊。\par
寄人·張泌\par
別夢依依到謝家，小廊回合曲闌斜。\par
多情只有春庭月，猶爲離人照落花。\par
題金陵渡·張祜\par
金陵津渡小山樓，一宿行人自可愁。\par
潮落夜江斜月裏，兩三星火是瓜州。\par
贈內人·張祜\par
禁門宮樹月痕過，媚眼惟看宿鷺窠。\par
斜拔玉釵燈影畔，剔開紅焰救飛蛾。\par
集靈臺·其一·張祜\par
日光斜照集靈臺，紅樹花迎曉露開。\par
昨夜上皇新授籙，太真含笑入簾來。\par
集靈臺·其二·張祜\par
虢國夫人承主恩，平明騎馬入宮門。\par
卻嫌脂粉污顏色，淡掃蛾眉朝至尊。\par
楓橋夜泊·張繼\par
月落烏啼霜滿天，江楓漁火對愁眠。\par
姑蘇城外寒山寺，夜半鐘聲到客船。\par
隴西行·陳陶\par
誓掃匈奴不顧身，五千貂錦喪胡塵。\par
可憐無定河邊骨，猶是春閨夢裏人！(春閨 一作：深閨)\par
雜詩·佚名\par
近寒食雨草萋萋，著麥苗風柳映堤。\par
等是有家歸未得，杜鵑休向耳邊啼。\par
回鄉偶書·其一·賀知章\par
少小離家老大回，鄉音無改鬢毛衰。\par
兒童相見不相識，笑問客從何處來。\par
瑤瑟怨·溫庭筠\par
冰簟銀牀夢不成，碧天如水夜雲輕。\par
雁聲遠過瀟湘去，十二樓中月自明。\par
月夜·劉方平\par
更深月色半人家，北斗闌干南鬥斜。\par
今夜偏知春氣暖，蟲聲新透綠窗紗。\par
春怨·劉方平\par
紗窗日落漸黃昏，金屋無人見淚痕。\par
寂寞空庭春欲晚，梨花滿地不開門。\par
烏衣巷·劉禹錫\par
朱雀橋邊野草花，烏衣巷口夕陽斜。\par
舊時王謝堂前燕，飛入尋常百姓家。\par
春詞·劉禹錫\par
新妝宜面下朱樓，深鎖春光一院愁。\par
行到中庭數花朵，蜻蜓飛上玉搔頭。\par
馬嵬坡·鄭畋\par
玄宗回馬楊妃死，雲雨雖亡日月新。\par
終是聖明天子事，景陽宮井又何人。\par
寒食·韓翃\par
春城無處不飛花，寒食東風御柳斜。\par
日暮漢宮傳蠟燭，輕煙散入五侯家。\par
已涼·韓偓\par
碧闌干外繡簾垂，猩血屏風畫折枝。\par
八尺龍鬚方錦褥，已涼天氣未寒時。\par
宮詞·顧況\par
長樂宮連上苑春，玉樓金殿豔歌新。\par
君門一入無由出，唯有宮鶯得見人。\par

\section{第3则}
送杜少府之任蜀州·王勃\par
城闕輔三秦，風煙望五津。\par
與君離別意，同是宦遊人。\par
海內存知己，天涯若比鄰。\par
無爲在歧路，兒女共沾巾。\par
送梓州李使君·王維\par
萬壑樹參天，千山響杜鵑。\par
山中一夜雨，樹杪百重泉。\par
漢女輸橦布，巴人訟芋田。\par
文翁翻教授，不敢倚先賢。\par
漢江臨眺·王維\par
楚塞三湘接，荊門九派通。\par
江流天地外，山色有無中。\par
郡邑浮前浦，波瀾動遠空。\par
襄陽好風日，留醉與山翁。\par
終南別業·王維\par
中歲頗好道，晚家南山陲。\par
興來每獨往，勝事空自知。\par
行到水窮處，坐看雲起時。\par
偶然值林叟，談笑無還期。\par
終南山·王維\par
太乙近天都，連山接海隅。\par
白雲回望合，青靄入看無。\par
分野中峯變，陰晴衆壑殊。\par
欲投人處宿，隔水問樵夫。\par
酬張少府·王維\par
晚年唯好靜，萬事不關心。自顧無長策，空知返舊林。\par
松風吹解帶，山月照彈琴。君問窮通理，漁歌入浦深。\par
過香積寺·王維\par
不知香積寺，數裏入雲峯。\par
古木無人徑，深山何處鐘。\par
泉聲咽危石，日色冷青松。\par
薄暮空潭曲，安禪製毒龍。\par
輞川閒居贈裴秀才迪·王維\par
寒山轉蒼翠，秋水日潺湲。\par
倚杖柴門外，臨風聽暮蟬。\par
渡頭餘落日，墟里上孤煙。\par
復值接輿醉，狂歌五柳前。\par
山居秋暝·王維\par
空山新雨後，天氣晚來秋。\par
明月鬆間照，清泉石上流。\par
竹喧歸浣女，蓮動下漁舟。\par
隨意春芳歇，王孫自可留。\par
歸嵩山作·王維\par
清川帶長薄，車馬去閒閒。\par
流水如有意，暮禽相與還。\par
荒城臨古渡，落日滿秋山。\par
迢遞嵩高下，歸來且閉關。\par
次北固山下·王灣\par
客路青山外，行舟綠水前。(青山外 一作：青山下)\par
潮平兩岸闊，風正一帆懸。\par
海日生殘夜，江春入舊年。\par
鄉書何處達？歸雁洛陽邊。\par
雲陽館與韓紳宿別·司空曙\par
故人江海別，幾度隔山川。\par
乍見翻疑夢，相悲各問年。\par
孤燈寒照雨，深竹暗浮煙。\par
更有明朝恨，離杯惜共傳。\par
喜外弟盧綸見宿·司空曙\par
靜夜四無鄰，荒居舊業貧。\par
雨中黃葉樹，燈下白頭人。\par
以我獨沉久，愧君相見頻。\par
平生自有分，況是蔡家親。\par
賊平後送人北歸·司空曙\par
世亂同南去，時清獨北還。\par
他鄉生白髮，舊國見青山。\par
曉月過殘壘，繁星宿故關。\par
寒禽與衰草，處處伴愁顏。\par
賦得古原草送別·白居易\par
離離原上草，一歲一枯榮。\par
野火燒不盡，春風吹又生。\par
遠芳侵古道，晴翠接荒城。\par
又送王孫去，萋萋滿別情。\par
題大庾嶺北驛·宋之問\par
陽月南飛雁，傳聞至此回。\par
我行殊未已，何日復歸來。\par
江靜潮初落，林昏瘴不開。\par
明朝望鄉處，應見隴頭梅。\par
寄左省杜拾遺·岑參\par
聯步趨丹陛，分曹限紫微。\par
曉隨天仗入，暮惹御香歸。\par
白髮悲花落，青雲羨鳥飛。\par
聖朝無闕事，自覺諫書稀。\par
聽蜀僧浚彈琴·李白\par
蜀僧抱綠綺，西下峨眉峯。\par
爲我一揮手，如聽萬壑鬆。\par
客心洗流水，餘響入霜鍾。\par
不覺碧山暮，秋雲暗幾重。\par
夜泊牛渚懷古·李白\par
牛渚西江夜，青天無片雲。\par
登舟望秋月，空憶謝將軍。\par
餘亦能高詠，斯人不可聞。\par
明朝掛帆席，楓葉落紛紛。(掛帆\par
一作：去)\par
贈孟浩然·李白\par
吾愛孟夫子，風流天下聞。\par
紅顏棄軒冕，白首臥鬆雲。\par
醉月頻中聖，迷花不事君。\par
高山安可仰，徒此揖清芬。\par
渡荊門送別·李白\par
渡遠荊門外，來從楚國遊。\par
山隨平野盡，江入大荒流。\par
月下飛天鏡，雲生結海樓。\par
仍憐故鄉水，萬里送行舟。\par
送友人·李白\par
青山橫北郭，白水繞東城。\par
此地一爲別，孤蓬萬里徵。\par
浮雲遊子意，落日故人情。\par
揮手自茲去，蕭蕭班馬鳴。\par
喜見外弟又言別·李益\par
十年離亂後，長大一相逢。\par
問姓驚初見，稱名憶舊容。\par
別來滄海事，語罷暮天鍾。\par
明日巴陵道，秋山又幾重。\par
涼思·李商隱\par
客去波平檻，蟬休露滿枝。永懷當此節，倚立自移時。\par
北斗兼春遠，南陵寓使遲。天涯占夢數，疑誤有新知。\par
北青蘿·李商隱\par
殘陽西入崦，茅屋訪孤僧。落葉人何在，寒雲路幾層。\par
獨敲初夜磬，閒倚一枝藤。世界微塵裏，吾寧愛與憎。\par
蟬·李商隱\par
本以高難飽，徒勞恨費聲。\par
五更疏欲斷，一樹碧無情。\par
薄宦梗猶泛，故園蕪已平。\par
煩君最相警，我亦舉家清。\par
風雨·李商隱\par
淒涼寶劍篇，羈泊欲窮年。\par
黃葉仍風雨，青樓自管絃。\par
新知遭薄俗，舊好隔良緣。\par
心斷新豐酒，銷愁鬥幾千。\par
落花·李商隱\par
高閣客竟去，小園花亂飛。參差連曲陌，迢遞送斜暉。\par
腸斷未忍掃，眼穿仍欲歸。芳心向春盡，所得是沾衣。\par
登岳陽樓·杜甫\par
昔聞洞庭水，今上岳陽樓。\par
吳楚東南坼，乾坤日夜浮。\par
親朋無一字，老病有孤舟。\par
戎馬關山北，憑軒涕泗流。\par
奉濟驛重送嚴公四韻·杜甫\par
遠送從此別，青山空復情。幾時杯重把，昨夜月同行。\par
列郡謳歌惜，三朝出入榮。江村獨歸處，寂寞養殘生。\par
別房太尉墓·杜甫\par
他鄉復行役，駐馬別孤墳。\par
近淚無干土，低空有斷雲。\par
對棋陪謝傅，把劍覓徐君。\par
唯見林花落，鶯啼送客聞。\par
旅夜書懷·杜甫\par
細草微風岸，危檣獨夜舟。\par
星垂平野闊，月涌大江流。\par
名豈文章著，官應老病休。\par
飄飄何所似，天地一沙鷗。\par
至德二載甫自京金光門·杜甫\par
此道昔歸順，西郊胡正繁。\par
至今殘破膽，應有未招魂。\par
近得歸京邑，移官豈至尊。\par
無才日衰老，駐馬望千門。\par
月夜憶舍弟·杜甫\par
戍鼓斷人行，邊秋一雁聲。(邊秋 一作:秋邊)\par
露從今夜白，月是故鄉明。\par
有弟皆分散，無家問死生。\par
寄書長不達，況乃未休兵。\par
天末懷李白·杜甫\par
涼風起天末，君子意如何。\par
鴻雁幾時到，江湖秋水多。\par
文章憎命達，魑魅喜人過。\par
應共冤魂語，投詩贈汨羅。\par
月夜·杜甫\par
今夜鄜州月，閨中只獨看。遙憐小兒女，未解憶長安。\par
香霧雲鬟溼，清輝玉臂寒。何時倚虛幌，雙照淚痕幹。\par
春望·杜甫\par
國破山河在，城春草木深。\par
感時花濺淚，恨別鳥驚心。\par
烽火連三月，家書抵萬金。\par
白頭搔更短，渾欲不勝簪。\par
春宿左省·杜甫\par
花隱掖垣暮，啾啾棲鳥過。\par
星臨萬戶動，月傍九霄多。\par
不寢聽金鑰，因風想玉珂。\par
明朝有封事，數問夜如何。\par
旅宿·杜牧\par
旅館無良伴，凝情自悄然。\par
寒燈思舊事，斷雁警愁眠。\par
遠夢歸侵曉，家書到隔年。\par
滄江好煙月，門系釣魚船。\par
春宮怨·杜荀鶴\par
早被嬋娟誤，欲妝臨鏡慵。承恩不在貌，教妾若爲容。\par
風暖鳥聲碎，日高花影重。年年越溪女，相憶採芙蓉。\par
和晉陵陸丞早春遊望·杜審言\par
獨有宦遊人，偏驚物候新。\par
雲霞出海曙，梅柳渡江春。\par
淑氣催黃鳥，晴光轉綠蘋。\par
忽聞歌古調，歸思欲沾巾。\par
雜詩三首·其三·沈佺期\par
聞道黃龍戍，頻年不解兵。\par
可憐閨裏月，長在漢家營。\par
少婦今春意，良人昨夜情。\par
誰能將旗鼓，一爲取龍城。\par
宿桐廬江寄廣陵舊遊·孟浩然\par
山暝聞猿愁，滄江急夜流。(聞 一作：聽)\par
風鳴兩岸葉，月照一孤舟。\par
建德非吾土，維揚憶舊遊。\par
還將兩行淚，遙寄海西頭。\par
留別王維·孟浩然\par
寂寂竟何待，朝朝空自歸。欲尋芳草去，惜與故人違。\par
當路誰相假，知音世所稀。只應守寂寞，還掩故園扉。\par
早寒有懷·孟浩然\par
木落雁南度，北風江上寒。\par
我家襄水曲，遙隔楚雲端。\par
鄉淚客中盡，孤帆天際看。(孤帆 一作：歸帆)\par
迷津欲有問，平海夕漫漫。\par
歲暮歸南山·孟浩然\par
北闕休上書，南山歸敝廬。不才明主棄，多病故人疏。\par
白髮催年老，青陽逼歲除。永懷愁不寐，鬆月夜窗虛。\par
過故人莊·孟浩然\par
故人具雞黍，邀我至田家。\par
綠樹村邊合，青山郭外斜。\par
開軒面場圃，把酒話桑麻。\par
待到重陽日，還來就菊花。\par
秦中寄遠上人·孟浩然\par
一丘常欲臥，三徑苦無資。\par
北土非吾願，東林懷我師。\par
黃金燃桂盡，壯志逐年衰。\par
日夕涼風至，聞蟬但益悲。\par
望洞庭湖贈張丞相·孟浩然\par
八月湖水平，涵虛混太清。\par
氣蒸雲夢澤，波撼岳陽城。\par
欲濟無舟楫，端居恥聖明。\par
坐觀垂釣者，徒有羨魚情。\par
與諸子登峴山·孟浩然\par
人事有代謝，往來成古今。\par
江山留勝蹟，我輩復登臨。\par
水落魚梁淺，天寒夢澤深。\par
羊公碑尚在，讀罷淚沾襟。(尚在 一作：字在)\par
宴梅道士山房·孟浩然\par
林臥愁春盡，開軒覽物華。\par
忽逢青鳥使，邀入赤松家。\par
金竈初開火，仙桃正發花。\par
童顏若可駐，何惜醉流霞。\par
章臺夜思·韋莊\par
清瑟怨遙夜，繞弦風雨哀。孤燈聞楚角，殘月下章臺。\par
芳草已雲暮，故人殊未來。鄉書不可寄，秋雁又南迴。\par
淮上喜會梁州故人·韋應物\par
江漢曾爲客，相逢每醉還。\par
浮雲一別後，流水十年間。\par
歡笑情如舊，蕭疏鬢已斑。\par
何因不歸去？淮上有秋山。\par
賦得暮雨送李曹·韋應物\par
楚江微雨裏，建業暮鍾時。\par
漠漠帆來重，冥冥鳥去遲。\par
海門深不見，浦樹遠含滋。\par
相送情無限，沾襟比散絲。\par
經鄒魯祭孔子而嘆之·李隆基\par
夫子何爲者，棲棲一代中。\par
地猶鄹氏邑，宅即魯王宮。\par
嘆鳳嗟身否，傷麟怨道窮。\par
今看兩楹奠，當與夢時同。\par
灞上秋居·馬戴\par
灞原風雨定，晚見雁行頻。\par
落葉他鄉樹，寒燈獨夜人。\par
空園白露滴，孤壁野僧鄰。\par
寄臥郊扉久，何年致此身。\par
楚江懷古三首·其一·馬戴\par
露氣寒光集，微陽下楚丘。猿啼洞庭樹，人在木蘭舟。\par
廣澤生明月，蒼山夾亂流。雲中君不見，竟夕自悲秋。\par
除夜有懷·崔塗\par
迢遞三巴路，羈危萬里身。\par
亂山殘雪夜，孤燭異鄉人。\par
漸與骨肉遠，轉於僮僕親。\par
那堪正飄泊，明日歲華新。\par
孤雁·崔塗\par
幾行歸塞盡，念爾獨何之。\par
暮雨相呼失，寒塘欲下遲。\par
渚雲低暗度，關月冷相隨。\par
未必逢矰繳，孤飛自可疑。\par
題破山寺後禪院·常建\par
清晨入古寺，初日照高林。\par
曲徑通幽處，禪房花木深。\par
山光悅鳥性，潭影空人心。\par
萬籟此俱寂，唯聞鐘磬音。\par
望月懷遠·張九齡\par
海上生明月，天涯共此時。\par
情人怨遙夜，竟夕起相思。\par
滅燭憐光滿，披衣覺露滋。\par
不堪盈手贈，還寢夢佳期。\par
書邊事·張喬\par
調角斷清秋，徵人倚戍樓。春風對青冢，白日落梁州。\par
大漢無兵阻，窮邊有客遊。蕃情似此水，長願向南流。\par
沒蕃故人·張籍\par
前年伐月支，城下沒全師。\par
蕃漢斷消息，死生長別離。\par
無人收廢帳，歸馬識殘旗。\par
欲祭疑君在，天涯哭此時。\par
秋日赴闕題潼關驛樓·許渾\par
紅葉晚蕭蕭，長亭酒一瓢。\par
殘雲歸太華，疏雨過中條。\par
樹色隨山迥，河聲入海遙。\par
帝鄉明日到，猶自夢漁樵。\par
早秋·許渾\par
遙夜泛清瑟，西風生翠蘿。\par
殘螢委玉露，早雁拂金河。\par
高樹曉還密，遠山晴更多。\par
淮南一葉下，自覺洞庭波。\par
送人東遊·溫庭筠\par
荒戍落黃葉，浩然離故關。\par
高風漢陽渡，初日郢門山。\par
江上幾人在，天涯孤棹還。\par
何當重相見，樽酒慰離顏。\par
尋陸鴻漸不遇·皎然\par
移家雖帶郭，野徑入桑麻。近種籬邊菊，秋來未著花。\par
扣門無犬吠，欲去問西家。報道山中去，歸時每日斜。\par
尋南溪常道士·劉長卿\par
一路經行處，莓苔見履痕。\par
白雲依靜渚，春草閉閒門。\par
過雨看鬆色，隨山到水源。\par
溪花與禪意，相對亦忘言。\par
新年作·劉長卿\par
鄉心新歲切，天畔獨潸然。\par
老至居人下，春歸在客先。\par
嶺猿同旦暮，江柳共風煙。\par
已似長沙傅，從今又幾年。\par
秋日登吳公臺上寺遠眺·劉長卿\par
古臺搖落後，秋日望鄉心。\par
野寺來人少，雲峯隔水深。\par
夕陽依舊壘，寒磬滿空林。\par
惆悵南朝事，長江獨至今。\par
送李中丞歸漢陽別業·劉長卿\par
流落徵南將，曾驅十萬師。\par
罷歸無舊業，老去戀明時。\par
獨立三邊靜，輕生一劍知。\par
茫茫江漢上，日暮欲何之。\par
餞別王十一南遊·劉長卿\par
望君煙水闊，揮手淚沾巾。\par
飛鳥沒何處，青山空向人。\par
長江一帆遠，落日五湖春。\par
誰見汀洲上，相思愁白蘋。\par
蜀先主廟·劉禹錫\par
天地英雄氣，千秋尚凜然。勢分三足鼎，業復五銖錢。\par
得相能開國，生兒不象賢。淒涼蜀故妓，來舞魏宮前。\par
闕題·劉昚虛\par
道由白雲盡，春與青溪長。\par
時有落花至，遠隨流水香。\par
閒門向山路，深柳讀書堂。\par
幽映每白日，清輝照衣裳。\par
送李端·盧綸\par
故關衰草遍，離別自堪悲。\par
路出寒雲外，人歸暮雪時。\par
少孤爲客早，多難識君遲。\par
掩淚空相向，風塵何處期。\par
送僧歸日本·錢起\par
上國隨緣住，來途若夢行。\par
浮天滄海遠，去世法舟輕。\par
水月通禪寂，魚龍聽梵聲。\par
惟憐一燈影，萬里眼中明。\par
谷口書齋寄楊補闕·錢起\par
泉壑帶茅茨，雲霞生薜帷。\par
竹憐新雨後，山愛夕陽時。\par
閒鷺棲常早，秋花落更遲。\par
家僮掃蘿徑，昨與故人期。\par
在獄詠蟬·駱賓王\par
西陸蟬聲唱，南冠客思深。(客思深 一作：客思侵)\par
不堪玄鬢影，來對白頭吟。(不堪 一作：那堪)\par
露重飛難進，風多響易沉。\par
無人信高潔，誰爲表予心？\par
江鄉故人偶集客舍·戴叔倫\par
天秋月又滿，城闕夜千重。\par
還作江南會，翻疑夢裏逢。\par
風枝驚暗鵲，露草覆寒蛩。\par
羈旅長堪醉，相留畏曉鍾。\par
酬程延秋夜即事見贈·韓翃\par
長簟迎風早，空城澹月華。\par
星河秋一雁，砧杵夜千家。\par
節候看應晚，心期臥亦賒。\par
向來吟秀句，不覺已鳴鴉。\par

\section{第4则}
遣悲懷三首·其一·元稹\par
謝公最小偏憐女，自嫁黔婁百事乖。\par
顧我無衣搜藎篋，泥他沽酒拔金釵。(藎篋 一作：畫篋)\par
野蔬充膳甘長藿，落葉添薪仰古槐。\par
今日俸錢過十萬，與君營奠復營齋。\par
遣悲懷三首·其二·元稹\par
昔日戲言身後事，今朝都到眼前來。\par
衣裳已施行看盡，針線猶存未忍開。\par
尚想舊情憐婢僕，也曾因夢送錢財。\par
誠知此恨人人有，貧賤夫妻百事哀。\par
遣悲懷三首·其三·元稹\par
閒坐悲君亦自悲，百年都是幾多時。\par
鄧攸無子尋知命，潘岳悼亡猶費詞。\par
同穴窅冥何所望，他生緣會更難期。\par
惟將終夜常開眼，報答平生未展眉。\par
贈郭給事·王維\par
洞門高閣靄餘輝，桃李陰陰柳絮飛。\par
禁裏疏鍾官舍晚，省中啼鳥吏人稀。\par
晨搖玉佩趨金殿，夕奉天書拜瑣闈。\par
強欲從君無那老，將因臥病解朝衣。\par
和賈至舍人早朝大明宮·王維\par
絳幘雞人送曉籌，尚衣方進翠雲裘。\par
九天閶闔開宮殿，萬國衣冠拜冕旒。\par
日色才臨仙掌動，香菸欲傍袞龍浮。\par
朝罷須裁五色詔，佩聲歸向鳳池頭。\par
奉和聖制從蓬萊向興慶閣道·王維\par
渭水自縈秦塞曲，黃山舊繞漢宮斜。\par
鑾輿迥出千門柳，閣道回看上苑花。\par
雲裏帝城雙鳳闕，雨中春樹萬人家。\par
爲乘陽氣行時令，不是宸遊玩物華。(玩 一作：重)\par
積雨輞川莊作·王維\par
積雨空林煙火遲，蒸藜炊黍餉東菑。\par
漠漠水田飛白鷺，陰陰夏木囀黃鸝。\par
山中習靜觀朝槿，松下清齋折露葵。\par
野老與人爭席罷，海鷗何事更相疑。\par
望月有感·白居易\par
時難年荒世業空，弟兄羈旅各西東。\par
田園寥落干戈後，骨肉流離道路中。\par
弔影分爲千里雁，辭根散作九秋蓬。\par
共看明月應垂淚，一夜鄉心五處同。\par
奉和中書舍人賈至早朝·岑參\par
雞鳴紫陌曙光寒，鶯囀皇州春色闌。\par
金闕曉鍾開萬戶，玉階仙仗擁千官。\par
花迎劍珮星初落，柳拂旌旗露未乾。\par
獨有鳳凰池上客，陽春一曲和皆難。\par
登金陵鳳凰臺·李白\par
鳳凰臺上鳳凰遊，鳳去臺空江自流。\par
吳宮花草埋幽徑，晉代衣冠成古丘。\par
三山半落青天外，二水中分白鷺洲。(二水 一作：一水)\par
總爲浮雲能蔽日，長安不見使人愁。\par
無題·重幃深下莫愁堂·李商隱\par
重幃深下莫愁堂，臥後清宵細細長。\par
神女生涯原是夢，小姑居處本無郎。\par
風波不信菱枝弱，月露誰教桂葉香。\par
直道相思了無益，未妨惆悵是清狂。\par
無題·鳳尾香羅薄幾重·李商隱\par
鳳尾香羅薄幾重，碧文圓頂夜深縫。\par
扇裁月魄羞難掩，車走雷聲語未通。\par
曾是寂寥金燼暗，斷無消息石榴紅。\par
斑騅只系垂楊岸，何處西南任好風。(任 一作：待)\par
無題·相見時難別亦難·李商隱\par
相見時難別亦難，東風無力百花殘。\par
春蠶到死絲方盡，蠟炬成灰淚始幹。\par
曉鏡但愁雲鬢改，夜吟應覺月光寒。\par
蓬山此去無多路，青鳥殷勤爲探看。(蓬山 一作：蓬萊)\par
無題·來是空言去絕蹤·李商隱\par
來是空言去絕蹤，月斜樓上五更鐘。\par
夢爲遠別啼難喚，書被催成墨未濃。\par
蠟照半籠金翡翠，麝薰微度繡芙蓉。\par
劉郎已恨蓬山遠，更隔蓬山一萬重！\par
無題·颯颯東風細雨來·李商隱\par
颯颯東風細雨來，芙蓉塘外有輕雷。\par
金蟾齧鎖燒香入，玉虎牽絲汲井回。\par
賈氏窺簾韓掾少，宓妃留枕魏王才。\par
春心莫共花爭發，一寸相思一寸灰！\par
無題·昨夜星辰昨夜風·李商隱\par
昨夜星辰昨夜風，畫樓西畔桂堂東。\par
身無綵鳳雙飛翼，心有靈犀一點通。\par
隔座送鉤春酒暖，分曹射覆蠟燈紅。\par
嗟餘聽鼓應官去，走馬蘭臺類轉蓬。\par
春雨·李商隱\par
悵臥新春白袷衣，白門寥落意多違。\par
紅樓隔雨相望冷，珠箔飄燈獨自歸。\par
遠路應悲春晼晚，殘霄猶得夢依稀。\par
玉璫緘札何由達，萬里雲羅一雁飛。\par
籌筆驛·李商隱\par
猿鳥猶疑畏簡書，風雲常爲護儲胥。\par
徒令上將揮神筆，終見降王走傳車。\par
管樂有才原不忝，關張無命欲何如？\par
他年錦裏經祠廟，梁父吟成恨有餘。\par
錦瑟·李商隱\par
錦瑟無端五十弦，一弦一柱思華年。\par
莊生曉夢迷蝴蝶，望帝春心託杜鵑。\par
滄海月明珠有淚，藍田日暖玉生煙。\par
此情可待成追憶，只是當時已惘然。\par
隋宮·李商隱\par
紫泉宮殿鎖煙霞，欲取蕪城作帝家。\par
玉璽不緣歸日角，錦帆應是到天涯。\par
於今腐草無螢火，終古垂楊有暮鴉。\par
地下若逢陳後主，豈宜重問後庭花。\par
送魏萬之京·李頎\par
朝聞遊子唱離歌，昨夜微霜初渡河。\par
鴻雁不堪愁裏聽，雲山況是客中過。\par
關城樹色催寒近，御苑砧聲向晚多。(\par
色 一作：曙)\par
莫見長安行樂處，空令歲月易蹉跎。\par
詠懷古蹟·其一·杜甫\par
支離東北風塵際，漂泊西南天地間。\par
三峽樓臺淹日月，五溪衣服共雲山。\par
羯胡事主終無賴，詞客哀時且未還。\par
庾信平生最蕭瑟，暮年詩賦動江關。\par
詠懷古蹟·其二·杜甫\par
搖落深知宋玉悲，風流儒雅亦吾師。\par
悵望千秋一灑淚，蕭條異代不同時。\par
江山故宅空文藻，雲雨荒臺豈夢思。\par
最是楚宮俱泯滅，舟人指點到今疑。\par
詠懷古蹟·其三·杜甫\par
羣山萬壑赴荊門，生長明妃尚有村。\par
一去紫臺連朔漠，獨留青冢向黃昏。\par
畫圖省識春風面，環佩空歸夜月魂。(夜月 一作：月夜 環\par
一作：環\par
)\par
千載琵琶作胡語，分明怨恨曲中論。\par
詠懷古蹟·其四·杜甫\par
蜀主窺吳幸三峽，崩年亦在永安宮。\par
翠華想像空山裏，玉殿虛無野寺中。\par
古廟杉鬆巢水鶴，歲時伏臘走村翁。\par
武侯祠堂常鄰近，一體君臣祭祀同。\par
詠懷古蹟·其五·杜甫\par
諸葛大名垂宇宙，宗臣遺像肅清高。\par
三分割據紆籌策，萬古雲霄一羽毛。\par
伯仲之間見伊呂，指揮若定失蕭曹。\par
運移漢祚終難復，志決身殲軍務勞。\par
宿府·杜甫\par
清秋幕府井梧寒，獨宿江城蠟炬殘。\par
永夜角聲悲自語，中天月色好誰看。\par
風塵荏苒音書絕，關塞蕭條行路難。\par
已忍伶俜十年事，強移棲息一枝安。\par
閣夜·杜甫\par
歲暮陰陽催短景，天涯霜雪霽寒宵。\par
五更鼓角聲悲壯，三峽星河影動搖。\par
野哭幾家聞戰伐，夷歌數處起漁樵。\par
臥龍躍馬終黃土，人事依依漫寂寥。\par
聞官軍收河南河北·杜甫\par
劍外忽傳收薊北，初聞涕淚滿衣裳。\par
卻看妻子愁何在，漫卷詩書喜欲狂。\par
白日放歌須縱酒，青春作伴好還鄉。\par
即從巴峽穿巫峽，便下襄陽向洛陽。\par
登高·杜甫\par
風急天高猿嘯哀，渚清沙白鳥飛回。\par
無邊落木蕭蕭下，不盡長江滾滾來。\par
萬里悲秋常作客，百年多病獨登臺。\par
艱難苦恨繁霜鬢，潦倒新停濁酒杯。\par
登樓·杜甫\par
花近高樓傷客心，萬方多難此登臨。\par
錦江春色來天地，玉壘浮雲變古今。\par
北極朝廷終不改，西山寇盜莫相侵。\par
可憐後主還祠廟，日暮聊爲樑甫吟。(樑甫 一作：樑父)\par
蜀相·杜甫\par
丞相祠堂何處尋，錦官城外柏森森。\par
映階碧草自春色，隔葉黃鸝空好音。\par
三顧頻煩天下計，兩朝開濟老臣心。(頻煩 一作：頻繁)\par
出師未捷身先死，長使英雄淚滿襟。\par
客至·杜甫\par
舍南舍北皆春水，但見羣鷗日日來。\par
花徑不曾緣客掃，蓬門今始爲君開。\par
盤飧市遠無兼味，樽酒家貧只舊醅。\par
肯與鄰翁相對飲，隔籬呼取盡餘杯。(餘 通：餘)\par
野望·杜甫\par
西山白雪三城戍，南浦清江萬里橋。\par
海內風塵諸弟隔，天涯涕淚一身遙。\par
惟將遲暮供多病，未有涓埃答聖朝。(惟 通：唯)\par
跨馬出郊時極目，不堪人事日蕭條。\par
登柳州城樓寄漳汀封連四州·柳宗元\par
城上高樓接大荒，海天愁思正茫茫。\par
驚風亂颭芙蓉水，密雨斜侵薜荔牆。\par
嶺樹重遮千里目，江流曲似九迴腸。\par
共來百越文身地，猶自音書滯一鄉。\par
春思·皇甫冉\par
鶯啼燕語報新年，馬邑龍堆路幾千。\par
家住層城臨漢苑，心隨明月到胡天。(層城 一作：秦城)\par
機中錦字論長恨，樓上花枝笑獨眠。\par
爲問元戎竇車騎，何時返旆勒燕然。\par
寄李儋元錫·韋應物\par
去年花裏逢君別，今日花開又一年。\par
世事茫茫難自料，春愁黯黯獨成眠。\par
身多疾病思田裏，邑有流亡愧俸錢。\par
聞道欲來相問訊，西樓望月幾回圓。\par
望薊門·祖詠\par
燕臺一望客心驚，笳鼓喧喧漢將營。(笳鼓 一作：簫鼓)\par
萬里寒光生積雪，三邊曙色動危旌。\par
沙場烽火侵胡月，海畔雲山擁薊城。\par
少小雖非投筆吏，論功還欲請長纓。\par
貧女·秦韜玉\par
蓬門未識綺羅香，擬託良媒益自傷。\par
誰愛風流高格調，共憐時世儉梳妝。\par
敢將十指誇針巧，不把雙眉鬥畫長。\par
苦恨年年壓金線，爲他人作嫁衣裳。\par
送李少府貶峽中王少府貶·高適\par
嗟君此別意何如，駐馬銜杯問謫居。\par
巫峽啼猿數行淚，衡陽歸雁幾封書。\par
青楓江上秋帆遠，白帝城邊古木疏。\par
聖代即今多雨露，暫時分手莫躊躇。\par
九日登望仙台呈劉明府·崔曙\par
漢文皇帝有高臺，此日登臨曙色開。\par
三晉雲山皆北向，二陵風雨自東來。\par
關門令尹誰能識，河上仙翁去不回。\par
且欲近尋彭澤宰，陶然共醉菊花杯。\par
登黃鶴樓·崔顥\par
昔人已乘黃鶴去，此地空餘黃鶴樓。\par
黃鶴一去不復返，白雲千載空悠悠。\par
晴川歷歷漢陽樹，芳草萋萋鸚鵡洲。\par
日暮鄉關何處是？煙波江上使人愁。\par
行經華陰·崔顥\par
岧嶢太華俯鹹京，天外三峯削不成。\par
武帝祠前雲欲散，仙人掌上雨初晴。\par
河山北枕秦關險，驛樹西連漢畤平。\par
借問路傍名利客，無如此處學長生。\par
利州南渡·溫庭筠\par
澹然空水對斜暉，曲島蒼茫接翠微。\par
波上馬嘶看棹去，柳邊人歇待船歸。\par
數叢沙草羣鷗散，萬頃江田一鷺飛。\par
誰解乘舟尋范蠡，五湖煙水獨忘機。\par
蘇武廟·溫庭筠\par
蘇武魂銷漢使前，古祠高樹兩茫然。\par
雲邊雁斷胡天月，隴上羊歸塞草煙。\par
回日樓臺非甲帳，去時冠劍是丁年。\par
茂陵不見封侯印，空向秋波哭逝川。\par
江州重別薛六柳·劉長卿\par
生涯豈料承優詔，世事空知學醉歌。\par
江上月明胡雁過，淮南木落楚山多。\par
寄身且喜滄洲近，顧影無如白髮何。\par
今日龍鍾人共棄，愧君猶遣慎風波。\par
長沙過賈誼宅·劉長卿\par
三年謫宦此棲遲，萬古惟留楚客悲。\par
秋草獨尋人去後，寒林空見日斜時。\par
漢文有道恩猶薄，湘水無情吊豈知。\par
寂寂江山搖落處，憐君何事到天涯。\par
自夏口至鸚鵡洲夕望岳陽·劉長卿\par
汀洲無浪復無煙，楚客相思益渺然。\par
漢口夕陽斜渡鳥，洞庭秋水遠連天。\par
孤城背嶺寒吹角，獨樹臨江夜泊船。(獨樹 一作：獨戍)\par
賈誼上書憂漢室，長沙謫去古今憐。\par
西塞山懷古·劉禹錫\par
王濬樓船下益州，金陵王氣黯然收。\par
千尋鐵鎖沉江底，一片降幡出石頭。\par
人世幾回傷往事，山形依舊枕寒流。\par
今逢四海爲家日，故壘蕭蕭蘆荻秋。\par
晚次鄂州·盧綸\par
雲開遠見漢陽城，猶是孤帆一日程。\par
估客晝眠知浪靜，舟人夜語覺潮生。\par
三湘愁鬢逢秋色，萬里歸心對月明。\par
舊業已隨征戰盡，更堪江上鼓鼙聲。\par
贈闕下裴舍人·錢起\par
二月黃鶯飛上林，春城紫禁曉陰陰。\par
長樂鐘聲花外盡，龍池柳色雨中深。\par
陽和不散窮途恨，霄漢長懷捧日心。\par
獻賦十年猶未遇，羞將白髮對華簪。\par
宮詞·薛逢\par
十二樓中盡曉妝，望仙樓上望君王。\par
鎖銜金獸連環冷，水滴銅龍晝漏長。\par
雲髻罷梳還對鏡，羅衣欲換更添香。\par
遙窺正殿簾開處，袍袴宮人掃御牀。\par
同題仙遊觀·韓翃\par
仙台初見五城樓，風物悽悽宿雨收。\par
山色遙連秦樹晚，砧聲近報漢宮秋。\par
疏鬆影落空壇靜，細草香閒小洞幽。\par
何用別尋方外去，人間亦自有丹丘。\par

\section{第5则}
賊退示官吏·元結\par
昔歲逢太平，山林二十年。\par
泉源在庭戶，洞壑當門前。\par
井稅有常期，日晏猶得眠。\par
忽然遭世變，數歲親戎旃。\par
今來典斯郡，山夷又紛然。\par
城小賊不屠，人貧傷可憐。\par
是以陷鄰境，此州獨見全。\par
使臣將王命，豈不如賊焉？\par
今彼征斂者，迫之如火煎。\par
誰能絕人命，以作時世賢！\par
思欲委符節，引竿自刺船。\par
將家就魚麥，歸老江湖邊。\par
同從弟銷南齋玩月憶山陰·王昌齡\par
高臥南齋時，開帷月初吐。\par
清輝澹水木，演漾在窗戶。\par
冉冉幾盈虛，澄澄變今古。\par
美人清江畔，是夜越吟苦。\par
千里共如何，微風吹蘭杜。\par
渭川田家·王維\par
斜陽照墟落，窮巷牛羊歸。(斜陽 一作：斜光)\par
野老念牧童，倚杖候荊扉。\par
雉雊麥苗秀，蠶眠桑葉稀。\par
田夫荷鋤至，相見語依依。\par
即此羨閒逸，悵然吟式微。\par
西施詠·王維\par
豔色天下重，西施寧久微。\par
朝爲越溪女，暮作吳宮妃。\par
賤日豈殊衆，貴來方悟稀。\par
邀人傅香粉，不自著羅衣。\par
君寵益嬌態，君憐無是非。\par
當時浣紗伴，莫得同車歸。\par
持謝鄰家子，效顰安可希。\par
送別·王維\par
下馬飲君酒，問君何所之？\par
君言不得意，歸臥南山陲。\par
但去莫復問，白雲無盡時。\par
送綦毋潛落第還鄉·王維\par
聖代無隱者，英靈盡來歸。\par
遂令東山客，不得顧采薇。\par
既至金門遠，孰雲吾道非。\par
江淮度寒食，京洛縫春衣。\par
置酒長安道，同心與我違。\par
行當浮桂棹，未幾拂荊扉。\par
遠樹帶行客，孤城當落暉。\par
吾謀適不用，勿謂知音稀。\par
青溪·王維\par
言入黃花川，每逐清溪水。\par
隨山將萬轉，趣途無百里。\par
聲喧亂石中，色靜深鬆裏。\par
漾漾泛菱荇，澄澄映葭葦。\par
我心素已閒，清川澹如此。\par
請留磐石上，垂釣將已矣。(磐石 一作：盤石)\par
與高適薛據同登慈恩寺·岑參\par
塔勢如涌出，孤高聳天宮。\par
登臨出世界，磴道盤虛空。\par
突兀壓神州，崢嶸如鬼工。\par
四角礙白日，七層摩蒼穹。\par
下窺指高鳥，俯聽聞驚風。\par
連山若波濤，奔湊似朝東。(湊 一作：走；似 一作：如)\par
青槐夾馳道，宮館何玲瓏。(館 一作：觀)\par
秋色從西來，蒼然滿關中。\par
五陵北原上，萬古青濛濛。\par
淨理了可悟，勝因夙所宗。\par
誓將掛冠去，覺道資無窮。\par
下終南山過斛斯山人宿置酒·李白\par
暮從碧山下，山月隨人歸。\par
卻顧所來徑，蒼蒼橫翠微。\par
相攜及田家，童稚開荊扉。\par
綠竹入幽徑，青蘿拂行衣。\par
歡言得所憩，美酒聊共揮。\par
長歌吟松風，曲盡河星稀。\par
我醉君復樂，陶然共忘機。\par
月下獨酌四首·其一·李白\par
花間一壺酒，獨酌無相親。\par
舉杯邀明月，對影成三人。\par
月既不解飲，影徒隨我身。\par
暫伴月將影，行樂須及春。\par
我歌月徘徊，我舞影零亂。\par
醒時同交歡，醉後各分散。(同交歡 一作：相交歡)\par
永結無情遊，相期邈雲漢。\par
春思·李白\par
燕草如碧絲，秦桑低綠枝。\par
當君懷歸日，是妾斷腸時。\par
春風不相識，何事入羅幃。\par
夢李白·其一·杜甫\par
死別已吞聲，生別常惻惻。江南瘴癘地，逐客無消息。\par
故人入我夢，明我長相憶。恐非平生魂，路遠不可測。\par
魂來楓葉青，魂返關塞黑。君今在羅網，何以有羽翼。\par
落月滿屋樑，猶疑照顏色。水深波浪闊，無使蛟龍得。\par
夢李白·其二·杜甫\par
浮雲終日行，遊子久不至。\par
三夜頻夢君，情親見君意。\par
告歸常侷促，苦道來不易。\par
江湖多風波，舟楫恐失墜。\par
出門搔白首，若負平生志。\par
冠蓋滿京華，斯人獨憔悴。\par
孰雲網恢恢，將老身反累。\par
千秋萬歲名，寂寞身後事。\par
望嶽·杜甫\par
岱宗夫如何？齊魯青未了。\par
造化鍾神秀，陰陽割昏曉。\par
蕩胸生曾雲，決眥入歸鳥。( 曾 通：層)\par
會當凌絕頂，一覽衆山小。\par
贈衛八處士·杜甫\par
人生不相見，動如參與商。\par
今夕復何夕，共此燈燭光！\par
少壯能幾時？鬢髮各已蒼！\par
訪舊半爲鬼，驚呼熱中腸。\par
佳人·杜甫\par
絕代有佳人，幽居在空谷。\par
自雲良家女，零落依草木。\par
關中昔喪亂，兄弟遭殺戮。\par
官高何足論，不得收骨肉。\par
世情惡衰歇，萬事隨轉燭。\par
夫婿輕薄兒，新人美如玉。\par
合昏尚知時，鴛鴦不獨宿。\par
但見新人笑，那聞舊人哭。\par
在山泉水清，出山泉水濁。\par
侍婢賣珠回，牽蘿補茅屋。\par
摘花不插發，採柏動盈掬。\par
天寒翠袖薄，日暮倚修竹。\par
秋登蘭山寄張五·孟浩然\par
北山白雲裏，隱者自怡悅。\par
相望試登高，心隨雁飛滅。(試 一作 始)\par
愁因薄暮起，興是清秋髮。\par
時見歸村人，沙行渡頭歇。\par
天邊樹若薺，江畔洲如月。\par
何當載酒來，共醉重陽節。\par
夏日南亭懷辛大·孟浩然\par
山光忽西落，池月漸東上。\par
散發乘夕涼，開軒臥閒敞。\par
荷風送香氣，竹露滴清響。\par
欲取鳴琴彈，恨無知音賞。\par
感此懷故人，中宵勞夢想。(\par
宵 一作：終)\par
宿業師山房待丁大不至·孟浩然\par
夕陽度西嶺，羣壑倏已暝。\par
鬆月生夜涼，風泉滿清聽。\par
樵人歸欲盡，煙鳥棲初定。\par
之子期宿來，孤琴候蘿徑。\par
尋西山隱者不遇·丘爲\par
絕頂一茅茨，直上三十里。\par
扣關無僮僕，窺室唯案几。\par
若非巾柴車，應是釣秋水。\par
差池不相見，黽勉空仰止。\par
草色新雨中，鬆聲晚窗裏。\par
及茲契幽絕，自足蕩心耳。\par
雖無賓主意，頗得清淨理。\par
興盡方下山，何必待之子。\par
晨詣超師院讀禪經·柳宗元\par
汲井漱寒齒，清心拂塵服。\par
閒持貝葉書，步出東齋讀。\par
真源了無取，妄跡世所逐。\par
遺言冀可冥，繕性何由熟。\par
道人庭宇靜，苔色連深竹。\par
日出霧露餘，青松如膏沐。\par
澹然離言說，悟悅心自足。\par
溪居·柳宗元\par
久爲簪組累，幸此南夷謫。\par
閒依農圃鄰，偶似山林客。\par
曉耕翻露草，夜榜響溪石。\par
來往不逢人，長歌楚天碧。\par
送楊氏女·韋應物\par
永日方慼慼，出行復悠悠。\par
女子今有行，大江溯輕舟。\par
爾輩苦無恃，撫念益慈柔。\par
幼爲長所育，兩別泣不休。\par
對此結中腸，義往難復留。\par
自小闕內訓，事姑貽我憂。\par
賴茲託令門，任恤庶無尤。\par
貧儉誠所尚，資從豈待周。\par
孝恭遵婦道，容止順其猷。\par
別離在今晨，見爾當何秋。\par
居閒始自遣，臨感忽難收。\par
歸來視幼女，零淚緣纓流。\par
長安遇馮著·韋應物\par
客從東方來，衣上灞陵雨。\par
問客何爲來，採山因買斧。\par
冥冥花正開，颺颺燕新乳。\par
昨別今已春，鬢絲生幾縷。\par
夕次盱眙縣·韋應物\par
落帆逗淮鎮，停舫臨孤驛。浩浩風起波，冥冥日沉夕。\par
人歸山郭暗，雁下蘆洲白。獨夜憶秦關，聽鍾未眠客。\par
東郊·韋應物\par
吏舍跼終年，出郊曠清曙。\par
楊柳散和風，青山澹吾慮。\par
依叢適自憩，緣澗還復去。\par
微雨靄芳原，春鳩鳴何處。\par
樂幽心屢止，遵事蹟猶遽。\par
終罷斯結廬，慕陶直可庶。\par
郡齋雨中與諸文士燕集·韋應物\par
兵衛森畫戟，宴寢凝清香。\par
海上風雨至，逍遙池閣涼。\par
煩痾近消散，嘉賓復滿堂。\par
自慚居處崇，未睹斯民康。\par
理會是非遣，性達形跡忘。\par
鮮肥屬時禁，蔬果幸見嘗。\par
俯飲一杯酒，仰聆金玉章。\par
神歡體自輕，意欲凌風翔。\par
吳中盛文史，羣彥今汪洋。\par
方知大藩地，豈曰財賦強。\par
初發揚子寄元大校書·韋應物\par
悽悽去親愛，泛泛入煙霧。\par
歸棹洛陽人，殘鍾廣陵樹。\par
今朝此爲別，何處還相遇。\par
世事波上舟，沿洄安得住。\par
寄全椒山中道士·韋應物\par
今朝郡齋冷，忽念山中客。\par
澗底束荊薪，歸來煮白石。\par
欲持一瓢酒，遠慰風雨夕。\par
落葉滿空山，何處尋行跡。\par
宿王昌齡隱居·常建\par
清溪深不測，隱處唯孤雲。\par
鬆際露微月，清光猶爲君。\par
茅亭宿花影，藥院滋苔紋。\par
餘亦謝時去，西山鸞鶴羣。\par
感遇·孤鴻海上來·張九齡\par
孤鴻海上來，池潢不敢顧。\par
側見雙翠鳥，巢在三珠樹。\par
矯矯珍木巔，得無金丸懼？\par
美服患人指，高明逼神惡。\par
今我遊冥冥，弋者何所慕！\par
感遇·蘭葉春葳蕤·張九齡\par
蘭葉春葳蕤，桂華秋皎潔。\par
欣欣此生意，自爾爲佳節。\par
誰知林棲者，聞風坐相悅。\par
草木有本心，何求美人折！\par
感遇·幽人歸獨臥·張九齡\par
幽人歸獨臥，滯慮洗孤清。\par
持此謝高鳥，因之傳遠情。\par
日夕懷空意，人誰感至精？\par
飛沈理自隔，何所慰吾誠?\par
感遇·江南有丹橘·張九齡\par
江南有丹橘，經冬猶綠林。\par
豈伊地氣暖？自有歲寒心。\par
可以薦嘉客，奈何阻重深。\par
運命惟所遇，循環不可尋。\par
徒言樹桃李，此木豈無陰？\par
春泛若耶溪·綦毋潛\par
幽意無斷絕，此去隨所偶。\par
晚風吹行舟，花路入溪口。\par
際夜轉西壑，隔山望南鬥。\par
潭煙飛溶溶，林月低向後。\par
生事且瀰漫，願爲持竿叟。\par

\section{第6则}
石魚湖上醉歌·元結\par
石魚湖，似洞庭，夏水欲滿君山青。\par
山爲樽，水爲沼，酒徒歷歷坐洲島。\par
長風連日作大浪，不能廢人運酒舫。\par
我持長瓢坐巴丘，酌飲四坐以散愁。\par
長恨歌·白居易\par
漢皇重色思傾國，御宇多年求不得。楊家有女初長成，養在深閨人未識。\par
天生麗質難自棄，一朝選在君王側。回眸一笑百媚生，六宮粉黛無顏色。\par
春寒賜浴華清池，溫泉水滑洗凝脂。侍兒扶起嬌無力，始是新承恩澤時。\par
雲鬢花顏金步搖，芙蓉帳暖度春宵。春宵苦短日高起，從此君王不早朝。\par
承歡侍宴無閒暇，春從春遊夜專夜。後宮佳麗三千人，三千寵愛在一身。\par
金屋妝成嬌侍夜，玉樓宴罷醉和春。姊妹弟兄皆列土，可憐光彩生門戶。\par
遂令天下父母心，不重生男重生女。驪宮高處入青雲，仙樂風飄處處聞。\par
緩歌慢舞凝絲竹，盡日君王看不足。漁陽鼙鼓動地來，驚破霓裳羽衣曲。\par
九重城闕煙塵生，千乘萬騎西南行。翠華搖搖行復止，西出都門百餘里。\par
六軍不發無奈何，宛轉蛾眉馬前死。花鈿委地無人收，翠翹金雀玉搔頭。\par
君王掩面救不得，回看血淚相和流。黃埃散漫風蕭索，雲棧縈紆登劍閣。\par
峨嵋山下少人行，旌旗無光日色薄。蜀江水碧蜀山青，聖主朝朝暮暮情。\par
行宮見月傷心色，夜雨聞鈴腸斷聲。天旋地轉回龍馭，到此躊躇不能去。\par
馬嵬坡下泥土中，不見玉顏空死處。君臣相顧盡沾衣，東望都門信馬歸。\par
歸來池苑皆依舊，太液芙蓉未央柳。芙蓉如面柳如眉，對此如何不淚垂。\par
春風桃李花開日，秋雨梧桐葉落時。西宮南內多秋草，落葉滿階紅不掃。(花開日 一作：花開夜；南內 一作：南苑)\par
梨園弟子白髮新，椒房阿監青娥老。夕殿螢飛思悄然，孤燈挑盡未成眠。\par
遲遲鐘鼓初長夜，耿耿星河欲曙天。鴛鴦瓦冷霜華重，翡翠衾寒誰與共。\par
悠悠生死別經年，魂魄不曾來入夢。臨邛道士鴻都客，能以精誠致魂魄。\par
爲感君王輾轉思，遂教方士殷勤覓。排空馭氣奔如電，昇天入地求之遍。\par
上窮碧落下黃泉，兩處茫茫皆不見。忽聞海上有仙山，山在虛無縹渺間。\par
樓閣玲瓏五雲起，其中綽約多仙子。中有一人字太真，雪膚花貌參差是。\par
金闕西廂叩玉扃，轉教小玉報雙成。聞道漢家天子使，九華帳裏夢魂驚。\par
攬衣推枕起徘徊，珠箔銀屏迤邐開。雲鬢半偏新睡覺，花冠不整下堂來。\par
風吹仙袂飄飄舉，猶似霓裳羽衣舞。玉容寂寞淚闌干，梨花一枝春帶雨。(闌 通：欄；飄飄 一作：飄颻)\par
含情凝睇謝君王，一別音容兩渺茫。昭陽殿裏恩愛絕，蓬萊宮中日月長。\par
回頭下望人寰處，不見長安見塵霧。惟將舊物表深情，鈿合金釵寄將去。\par
釵留一股合一扇，釵擘黃金合分鈿。但教心似金鈿堅，天上人間會相見。\par
臨別殷勤重寄詞，詞中有誓兩心知。七月七日長生殿，夜半無人私語時。\par
在天願作比翼鳥，在地願爲連理枝。天長地久有時盡，此恨綿綿無絕期。\par
琵琶行·白居易\par
潯陽江頭夜送客，楓葉荻花秋瑟瑟。\par
主人下馬客在船，舉酒欲飲無管絃。\par
醉不成歡慘將別，別時茫茫江浸月。\par
忽聞水上琵琶聲，主人忘歸客不發。\par
尋聲暗問彈者誰？琵琶聲停欲語遲。\par
移船相近邀相見，添酒回燈重開宴。\par
千呼萬喚始出來，猶抱琵琶半遮面。\par
轉軸撥絃三兩聲，未成曲調先有情。\par
弦弦掩抑聲聲思，似訴平生不得志。（不得志 一作：意）\par
低眉信手續續彈，說盡心中無限事。\par
輕攏慢捻抹復挑，初爲《霓裳》後《六幺》（六幺 一作：綠腰）。\par
大弦嘈嘈如急雨，小弦切切如私語。\par
嘈嘈切切錯雜彈，大珠小珠落玉盤。\par
間關鶯語花底滑，幽咽泉流冰下難。\par
冰泉冷澀弦凝絕，凝絕不通聲暫歇。(暫歇 一作：漸歇)\par
別有幽愁暗恨生，此時無聲勝有聲。\par
銀瓶乍破水漿迸，鐵騎突出刀槍鳴。\par
曲終收撥當心畫，四弦一聲如裂帛。\par
東船西舫悄無言，唯見江心秋月白。\par
沉吟放撥插弦中，整頓衣裳起斂容。\par
自言本是京城女，家在蝦蟆陵下住。\par
十三學得琵琶成，名屬教坊第一部。\par
曲罷曾教善才服，妝成每被秋娘妒。\par
五陵年少爭纏頭，一曲紅綃不知數。\par
鈿頭銀篦擊節碎，血色羅裙翻酒污。(銀篦 一作：雲)\par
今年歡笑復明年，秋月春風等閒度。\par
弟走從軍阿姨死，暮去朝來顏色故。\par
門前冷落鞍馬稀，老大嫁作商人婦。\par
商人重利輕別離，前月浮樑買茶去。\par
去來江口守空船，繞船月明江水寒。\par
夜深忽夢少年事，夢啼妝淚紅闌干。\par
我聞琵琶已嘆息，又聞此語重唧唧。\par
同是天涯淪落人，相逢何必曾相識！\par
我從去年辭帝京，謫居臥病潯陽城。\par
潯陽地僻無音樂，終歲不聞絲竹聲。\par
住近湓江地低溼，黃蘆苦竹繞宅生。\par
其間旦暮聞何物？杜鵑啼血猿哀鳴。\par
春江花朝秋月夜，往往取酒還獨傾。\par
豈無山歌與村笛？嘔啞嘲哳難爲聽。\par
今夜聞君琵琶語，如聽仙樂耳暫明。\par
莫辭更坐彈一曲，爲君翻作《琵琶行》。\par
感我此言良久立，卻坐促弦弦轉急。\par
悽悽不似向前聲，滿座重聞皆掩泣。\par
座中泣下誰最多？江州司馬青衫溼。\par
走馬川行奉送封大夫出師·岑參\par
君不見走馬川行雪海邊，平沙莽莽黃入天。\par
輪臺九月風夜吼，一川碎石大如鬥，隨風滿地石亂走。\par
匈奴草黃馬正肥，金山西見煙塵飛，漢家大將西出師。\par
將軍金甲夜不脫，半夜軍行戈相撥，風頭如刀面如割。\par
馬毛帶雪汗氣蒸，五花連錢旋作冰，幕中草檄硯水凝。\par
虜騎聞之應膽懾，料知短兵不敢接，車師西門佇獻捷。\par
輪臺歌奉送封大夫出師·岑參\par
輪臺城頭夜吹角，輪臺城北旄頭落。\par
羽書昨夜過渠黎，單于已在金山西。\par
戍樓西望煙塵黑，漢軍屯在輪臺北。\par
上將擁旄西出徵，平明吹笛大軍行。\par
四邊伐鼓雪海涌，三軍大呼陰山動。\par
虜塞兵氣連雲屯，戰場白骨纏草根。\par
劍河風急雪片闊，沙口石凍馬蹄脫。(雪片闊 一作：雲片闊)\par
亞相勤王甘苦辛，誓將報主靜邊塵。\par
古來青史誰不見，今見功名勝古人。\par
白雪歌送武判官歸京·岑參\par
北風捲地白草折，胡天八月即飛雪。\par
忽如一夜春風來，千樹萬樹梨花開。\par
散入珠簾溼羅幕，狐裘不暖錦衾薄。\par
將軍角弓不得控，都護鐵衣冷難着。(難着 一作：猶著)\par
瀚海闌干百丈冰，愁雲慘淡萬里凝。\par
中軍置酒飲歸客，胡琴琵琶與羌笛。\par
紛紛暮雪下轅門，風掣紅旗凍不翻。\par
輪臺東門送君去，去時雪滿天山路。\par
山迴路轉不見君，雪上空留馬行處。\par
宣州謝脁樓餞別校書叔雲·李白\par
棄我去者，昨日之日不可留；\par
亂我心者，今日之日多煩憂。\par
長風萬里送秋雁，對此可以酣高樓。\par
蓬萊文章建安骨，中間小謝又清發。\par
俱懷逸興壯思飛，欲上青天攬明月。(覽 通：攬；明月 一作：日月)\par
抽刀斷水水更流，舉杯消愁愁更愁。\par
人生在世不稱意，明朝散發弄扁舟。\par
廬山謠寄盧侍御虛舟·李白\par
我本楚狂人，鳳歌笑孔丘。\par
手持綠玉杖，朝別黃鶴樓。\par
五嶽尋仙不辭遠，一生好入名山遊。\par
廬山秀出南鬥傍，屏風九疊雲錦張。\par
影落明湖青黛光，金闕前開二峯長，銀河倒掛三石樑。\par
香爐瀑布遙相望，回崖沓嶂凌蒼蒼。\par
翠影紅霞映朝日，鳥飛不到吳天長。\par
登高壯觀天地間，大江茫茫去不還。\par
黃雲萬里動風色，白波九道流雪山。\par
好爲廬山謠，興因廬山發。\par
閒窺石鏡清我心，謝公行處蒼苔沒。\par
早服還丹無世情，琴心三疊道初成。\par
遙見仙人彩雲裏，手把芙蓉朝玉京。\par
先期汗漫九垓上，願接盧敖遊太清。\par
夢遊天姥吟留別·李白\par
海客談瀛洲，煙濤微茫信難求；\par
越人語天姥，雲霞明滅或可睹。\par
天姥連天向天橫，勢拔五嶽掩赤城。\par
天台四萬八千丈，對此欲倒東南傾。(四萬 一作：一萬)\par
我欲因之夢吳越，一夜飛度鏡湖月。(度 通：渡)\par
湖月照我影，送我至剡溪。\par
謝公宿處今尚在，淥水盪漾清猿啼。\par
腳著謝公屐，身登青雲梯。\par
半壁見海日，空中聞天雞。\par
千巖萬轉路不定，迷花倚石忽已暝。\par
熊咆龍吟殷巖泉，慄深林兮驚層巔。\par
雲青青兮欲雨，水澹澹兮生煙。\par
列缺霹靂，丘巒崩摧。\par
洞天石扉，訇然中開。\par
青冥浩蕩不見底，日月照耀金銀臺。\par
霓爲衣兮風爲馬，雲之君兮紛紛而來下。\par
虎鼓瑟兮鸞回車，仙之人兮列如麻。\par
忽魂悸以魄動，恍驚起而長嗟。\par
惟覺時之枕蓆，失向來之煙霞。\par
世間行樂亦如此，古來萬事東流水。\par
別君去兮何時還？且放白鹿青崖間，須行即騎訪名山。\par
安能摧眉折腰事權貴，使我不得開心顏!\par
金陵酒肆留別·李白\par
風吹柳花滿店香，吳姬壓酒喚客嘗。(勸客 一作：喚客)\par
金陵子弟來相送，欲行不行各盡觴。\par
請君試問東流水，別意與之誰短長。\par
韓碑·李商隱\par
元和天子神武姿，彼何人哉軒與羲。\par
誓將上雪列聖恥，坐法宮中朝四夷。\par
淮西有賊五十載，封狼生貙貙生羆。\par
不據山河據平地，長戈利矛日可麾。\par
帝得聖相相曰度，賊斫不死神扶持。\par
腰懸相印作都統，陰風慘澹天王旗。\par
愬武古通作牙爪，儀曹外郎載筆隨。\par
行軍司馬智且勇，十四萬衆猶虎貔。\par
入蔡縛賊獻太廟，功無與讓恩不訾。\par
帝曰汝度功第一，汝從事愈宜爲辭。\par
愈拜稽首蹈且舞，金石刻畫臣能爲。\par
古者世稱大手筆，此事不繫於職司。\par
當仁自古有不讓，言訖屢頷天子頤。\par
公退齋戒坐小閣，濡染大筆何淋漓。\par
點竄堯典舜典字，塗改清廟生民詩。\par
文成破體書在紙，清晨再拜鋪丹墀。\par
表曰臣愈昧死上，詠神聖功書之碑。\par
碑高三丈字如鬥，負以靈鰲蟠以螭。\par
句奇語重喻者少，讒之天子言其私。\par
長繩百尺拽碑倒，粗砂大石相磨治。\par
公之斯文若元氣，先時已入人肝脾。\par
湯盤孔鼎有述作，今無其器存其辭。\par
嗚呼聖王及聖相，相與烜赫流淳熙。\par
公之斯文不示後，曷與三五相攀追。\par
願書萬本誦萬遍，口角流沫右手胝。\par
傳之七十有二代，以爲封禪玉檢明堂基。\par
聽董大彈胡笳聲兼寄語·李頎\par
蔡女昔造胡笳聲，一彈一十有八拍。\par
胡人落淚沾邊草，漢使斷腸對歸客。\par
古戍蒼蒼烽火寒，大荒沉沉飛雪白。\par
先拂商弦后角羽，四郊秋葉驚摵摵。\par
董夫子，通神明，深山竊聽來妖精。\par
言遲更速皆應手，將往復旋如有情。\par
空山百鳥散還合，萬里浮雲陰且晴。\par
嘶酸雛雁失羣夜，斷絕胡兒戀母聲。\par
川爲靜其波，鳥亦罷其鳴。\par
烏孫部落家鄉遠，邏娑沙塵哀怨生。\par
幽音變調忽飄灑，長風吹林雨墮瓦。\par
迸泉颯颯飛木末，野鹿呦呦走堂下。\par
長安城連東掖垣，鳳凰池對青瑣門。\par
高才脫略名與利，日夕望君抱琴至。\par
聽安萬善吹觱篥歌·李頎\par
南山截竹爲觱篥，此樂本自龜茲出。\par
流傳漢地曲轉奇，涼州胡人爲我吹。\par
傍鄰聞者多嘆息，遠客思鄉皆淚垂。\par
世人解聽不解賞，長飆風中自來往。\par
枯桑老柏寒颼飀，九雛鳴鳳亂啾啾。\par
龍吟虎嘯一時發，萬籟百泉相與秋。\par
忽然更作漁陽摻，黃雲蕭條白日暗。\par
變調如聞楊柳春，上林繁花照眼新。\par
歲夜高堂列明燭，美酒一杯聲一曲。\par
古意·李頎\par
男兒事長征，少小幽燕客。\par
賭勝馬蹄下，由來輕七尺。\par
殺人莫敢前，須如蝟毛磔。\par
黃雲隴底白雲飛，未得報恩不能歸。\par
遼東小婦年十五，慣彈琵琶解歌舞。\par
今爲羌笛出塞聲，使我三軍淚如雨。\par
送陳章甫·李頎\par
四月南風大麥黃，棗花未落桐陰長。\par
青山朝別暮還見，嘶馬出門思舊鄉。\par
陳侯立身何坦蕩，虯鬚虎眉仍大顙。\par
腹中貯書一萬卷，不肯低頭在草莽。\par
東門酤酒飲我曹，心輕萬事如鴻毛。\par
醉臥不知白日暮，有時空望孤雲高。\par
長河浪頭連天黑，津口停舟渡不得。\par
鄭國遊人未及家，洛陽行子空嘆息。\par
聞道故林相識多，罷官昨日今如何。\par
琴歌·李頎\par
主人有酒歡今夕，請奏鳴琴廣陵客。\par
月照城頭烏半飛，霜悽萬樹風入衣。(萬樹 一作：萬木)\par
銅爐華燭燭增輝，初彈淥水後楚妃。\par
一聲已動物皆靜，四座無言星欲稀。\par
清淮奉使千餘里，敢告雲山從此始。\par
古柏行·杜甫\par
孔明廟前有老柏，柯如青銅根如石。\par
霜皮溜雨四十圍，黛色參天二千尺。\par
君臣已與時際會，樹木猶爲人愛惜。\par
雲來氣接巫峽長，月出寒通雪山白。\par
憶昨路繞錦亭東，先主武侯同閟宮。\par
崔嵬枝幹郊原古，窈窕丹青戶牖空。\par
落落盤踞雖得地，冥冥孤高多烈風。\par
扶持自是神明力，正直原因造化工。\par
大廈如傾要樑棟，萬牛回首丘山重。\par
不露文章世已驚，未辭翦伐誰能送？\par
苦心豈免容螻蟻，香葉終經宿鸞鳳。\par
志士幽人莫怨嗟：古來材大難爲用。\par
觀公孫大娘弟子舞劍器行·杜甫\par
昔有佳人公孫氏，一舞劍器動四方。\par
觀者如山色沮喪，天地爲之久低昂。\par
霍如羿射九日落，矯如羣帝驂龍翔。\par
來如雷霆收震怒，罷如江海凝清光。\par
絳脣珠袖兩寂寞，晚有弟子傳芬芳。\par
臨潁美人在白帝，妙舞此曲神揚揚。\par
與餘問答既有以，感時撫事增惋傷。\par
先帝侍女八千人，公孫劍器初第一。\par
五十年間似反掌，風塵澒洞昏王室。\par
梨園弟子散如煙，女樂餘姿映寒日。\par
金粟堆南木已拱，瞿唐石城草蕭瑟。\par
玳筵急管曲復終，樂極哀來月東出。\par
老夫不知其所往，足繭荒山轉愁疾。\par
韋諷錄事宅觀曹將軍畫馬·杜甫\par
國初已來畫鞍馬，神妙獨數江都王。\par
將軍得名三十載，人間又見真乘黃。\par
曾貌先帝照夜白，龍池十日飛霹靂。\par
內府殷紅瑪瑙盤，婕妤傳詔才人索。\par
盤賜將軍拜舞歸，輕紈細綺相追飛。\par
貴戚權門得筆跡，始覺屏障生光輝。\par
昔日太宗拳毛騧，近時郭家獅子花。\par
今之新圖有二馬，復令識者久嘆嗟。\par
此皆騎戰一敵萬，縞素漠漠開風沙。\par
其餘七匹亦殊絕，迥若寒空動煙雪。\par
霜蹄蹴踏長楸間，馬官廝養森成列。\par
可憐九馬爭神駿，顧視清高氣深穩。\par
借問苦心愛者誰，後有韋諷前支遁。\par
憶昔巡幸新豐宮，翠華拂天來向東。\par
騰驤磊落三萬匹，皆與此圖筋骨同。\par
自從獻寶朝河宗，無復射蛟江水中。\par
君不見金粟堆前松柏裏，龍媒去盡鳥呼風。\par
丹青引贈曹霸將軍·杜甫\par
將軍魏武之子孫，於今爲庶爲清門。\par
英雄割據雖已矣，文采風流今尚存。\par
學書初學衛夫人，但恨無過王右軍。\par
丹青不知老將至，富貴於我如浮雲。\par
開元之中常引見，承恩數上南薰殿。\par
凌煙功臣少顏色，將軍下筆開生面。\par
良相頭上進賢冠，猛將腰間大羽箭。\par
褒公鄂公毛髮動，英姿颯爽來酣戰。\par
先帝御馬五花驄，畫工如山貌不同。\par
是日牽來赤墀下，迥立閶闔生長風。\par
詔謂將軍拂絹素，意匠慘澹經營中。\par
斯須九重真龍出，一洗萬古凡馬空。\par
玉花卻在御榻上，榻上庭前屹相向。\par
至尊含笑催賜金，圉人太僕皆惆悵。\par
弟子韓幹早入室，亦能畫馬窮殊相。\par
幹惟畫肉不畫骨，忍使驊騮氣凋喪。\par
將軍畫善蓋有神，必逢佳士亦寫真。\par
即今漂泊干戈際，屢貌尋常行路人。\par
途窮反遭俗眼白，世上未有如公貧。\par
但看古來盛名下，終日坎壈纏其身。\par
寄韓諫議·杜甫\par
今我不樂思岳陽，身欲奮飛病在牀。\par
美人娟娟隔秋水，濯足洞庭望八荒。\par
鴻飛冥冥日月白，青楓葉赤天雨霜。\par
玉京羣帝集北斗，或騎麒麟翳鳳凰。\par
芙蓉旌旗煙霧落，影動倒景搖瀟湘。\par
星宮之君醉瓊漿，羽人稀少不在旁。\par
似聞昨者赤松子，恐是漢代韓張良。\par
昔隨劉氏定長安，帷幄未改神慘傷。\par
國家成敗吾豈敢，色難腥腐餐楓香。\par
周南留滯古所惜，南極老人應壽昌。\par
美人胡爲隔秋水，焉得置之貢玉堂。\par
夜歸鹿門山歌·孟浩然\par
山寺鐘鳴晝已昏，漁梁渡頭爭渡喧。\par
人隨沙岸向江村，餘亦乘舟歸鹿門。\par
鹿門月照開煙樹，忽到龐公棲隱處。\par
巖扉鬆徑長寂寥，惟有幽人自來去。\par
漁翁·柳宗元\par
漁翁夜傍西巖宿，曉汲清湘燃楚竹。\par
煙銷日出不見人，欸乃一聲山水綠。\par
回看天際下中流，巖上無心雲相逐。\par
登幽州臺歌·陳子昂\par
前不見古人，後不見來者。\par
念天地之悠悠，獨愴然而涕下。\par
石鼓歌·韓愈\par
張生手持石鼓文，勸我試作石鼓歌。\par
少陵無人謫仙死，才薄將奈石鼓何。\par
周綱凌遲四海沸，宣王憤起揮天戈。\par
大開明堂受朝賀，諸侯劍佩鳴相磨。\par
蒐於岐陽騁雄俊，萬里禽獸皆遮羅。\par
鐫功勒成告萬世，鑿石作鼓隳嵯峨。\par
從臣才藝鹹第一，揀選撰刻留山阿。\par
雨淋日灸野火燎，鬼物守護煩撝呵。\par
公從何處得紙本，毫髮盡備無差訛。\par
辭嚴義密讀難曉，字體不類隸與蝌。\par
年深豈免有缺畫，快劍斫斷生蛟鼉。\par
鸞翔鳳翥衆仙下，珊瑚碧樹交枝柯。\par
金繩鐵索鎖鈕壯，古鼎躍水龍騰梭。\par
陋儒編詩不收入，二雅褊迫無委蛇。\par
孔子西行不到秦，掎摭星宿遺羲娥。\par
嗟餘好古生苦晚，對此涕淚雙滂沱。\par
憶昔初蒙博士徵，其年始改稱元和。\par
故人從軍在右輔，爲我度量掘臼科。\par
濯冠沐浴告祭酒，如此至寶存豈多。\par
氈包席裹可立致，十鼓只載數駱駝。\par
薦諸太廟比郜鼎，光價豈止百倍過。\par
聖恩若許留太學，諸生講解得切磋。\par
觀經鴻都尚填咽，坐見舉國來奔波。\par
剜苔剔蘚露節角，安置妥帖平不頗。\par
大廈深檐與蓋覆，經歷久遠期無佗。\par
中朝大官老於事，詎肯感激徒媕婀。\par
牧童敲火牛礪角，誰復著手爲摩挲。\par
日銷月鑠就埋沒，六年西顧空吟哦。\par
羲之俗書趁姿媚，數紙尚可博白鵝。\par
繼周八代爭戰罷，無人收拾理則那。\par
方今太平日無事，柄任儒術崇丘軻。\par
安能以此尚論列，願借辯口如懸河。\par
石鼓之歌止於此，嗚呼吾意其蹉跎。\par
山石·韓愈\par
山石犖确行徑微，黃昏到寺蝙蝠飛。\par
升堂坐階新雨足，芭蕉葉大梔子肥。\par
僧言古壁佛畫好，以火來照所見稀。\par
鋪牀拂席置羹飯，疏糲亦足飽我飢。\par
夜深靜臥百蟲絕，清月出嶺光入扉。\par
天明獨去無道路，出入高下窮煙霏。\par
山紅澗碧紛爛漫，時見鬆櫪皆十圍。\par
當流赤足踏澗石，水聲激激風吹衣。\par
人生如此自可樂，豈必局束爲人鞿？(鞿 一作：靰)\par
嗟哉吾黨二三子，安得至老不更歸。\par
八月十五夜贈張功曹·韓愈\par
纖雲四卷天無河，清風吹空月舒波。\par
沙平水息聲影絕，一杯相屬君當歌。\par
君歌聲酸辭且苦，不能聽終淚如雨。\par
洞庭連天九疑高，蛟龍出沒猩鼯號。\par
十生九死到官所，幽居默默如藏逃。\par
下牀畏蛇食畏藥，海氣溼蟄薰腥臊。\par
昨者州前捶大鼓，嗣皇繼聖登夔皋。\par
赦書一日行萬里，罪從大辟皆除死。\par
遷者追回流者還，滌瑕盪垢清朝班。\par
州家申名使家抑，坎軻只得移荊蠻。\par
判司卑官不堪說，未免捶楚塵埃間。\par
同時輩流多上道，天路幽險難追攀。\par
君歌且休聽我歌，我歌今與君殊科。\par
一年明月今宵多，人生由命非由他。\par
有酒不飲奈明何。\par
謁衡岳廟遂宿嶽寺題門樓·韓愈\par
五嶽祭秩皆三公，四方環鎮嵩當中。\par
火維地荒足妖怪，天假神柄專其雄。\par
噴雲泄霧藏半腹，雖有絕頂誰能窮？\par
我來正逢秋雨節，陰氣晦昧無清風。\par
潛心默禱若有應，豈非正直能感通！\par
須臾靜掃衆峯出，仰見突兀撐青空。\par
紫蓋連延接天柱，石廩騰擲堆祝融。\par
森然魄動下馬拜，松柏一徑趨靈宮。\par
粉牆丹柱動光彩，鬼物圖畫填青紅。\par
升階傴僂薦脯酒，欲以菲薄明其衷。\par
廟令老人識神意，睢盱偵伺能鞠躬。\par
手持杯珓導我擲，雲此最吉餘難同。\par
竄逐蠻荒幸不死，衣食才足甘長終。\par
侯王將相望久絕，神縱慾福難爲功。\par
夜投佛寺上高閣，星月掩映雲曈曨。\par
猿鳴鐘動不知曙，杲杲寒日生於東。\par

\section{第7则}
涼州詞·王之渙\par
黃河遠上白雲間，一片孤城萬仞山。\par
羌笛何須怨楊柳，春風不度玉門關。\par
出塞·秦時明月漢時關·王昌齡\par
秦時明月漢時關，萬里長征人未還。\par
但使龍城飛將在，不教胡馬度陰山。\par
塞上曲·蟬鳴空桑林·王昌齡\par
蟬鳴空桑林，八月蕭關道。\par
出塞入塞寒，處處黃蘆草。\par
從來幽並客，皆共塵沙老。\par
莫學遊俠兒，矜誇紫騮好。\par
塞下曲·飲馬渡秋水·王昌齡\par
飲馬渡秋水，水寒風似刀。\par
平沙日未沒，黯黯見臨洮。\par
昔日長城戰，鹹言意氣高。\par
黃塵足今古，白骨亂蓬蒿。\par
長信怨·王昌齡\par
金井梧桐秋葉黃，珠簾不卷夜來霜。\par
熏籠玉枕無顏色，臥聽南宮清漏長。\par
高殿秋砧響夜闌，霜深猶憶御衣寒。\par
銀燈青瑣裁縫歇，還向金城明主看。\par
奉帚平明金殿開，暫將團扇共徘徊。\par
玉顏不及寒鴉色，猶帶昭陽日影來。\par
真成薄命久尋思，夢見君王覺後疑。\par
火照西宮知夜飲，分明覆道奉恩時。\par
長信宮中秋月明，昭陽殿下搗衣聲。\par
白露堂中細草跡，紅羅帳裏不勝情。\par
渭城曲·王維\par
渭城朝雨浥輕塵，客舍青青柳色新。(一作：客舍依依楊柳春)\par
勸君更盡一杯酒，西出陽關無故人。\par
秋夜曲·王維\par
桂魄初生秋露微，輕羅已薄未更衣。\par
銀箏夜久殷勤弄，心怯空房不忍歸。\par
洛陽女兒行·王維\par
洛陽女兒對門居，纔可顏容十五餘。\par
良人玉勒乘驄馬，侍女金盤膾鯉魚。\par
畫閣朱樓盡相望，紅桃綠柳垂檐向。\par
羅帷送上七香車，寶扇迎歸九華帳。\par
狂夫富貴在青春，意氣驕奢劇季倫。\par
自憐碧玉親教舞，不惜珊瑚持與人。\par
春窗曙滅九微火，九微片片飛花瑣。\par
戲罷曾無理曲時，妝成祗是薰香坐。\par
城中相識盡繁華，日夜經過趙李家。\par
誰憐越女顏如玉，貧賤江頭自浣紗。\par
老將行·王維\par
少年十五二十時，步行奪得胡馬騎。\par
射殺山中白額虎，肯數鄴下黃鬚兒！\par
一身轉戰三千里，一劍曾當百萬師。\par
漢兵奮迅如霹靂，虜騎崩騰畏蒺藜。\par
衛青不敗由天幸，李廣無功緣數奇。\par
自從棄置便衰朽，世事蹉跎成白首。\par
昔時飛箭無全目，今日垂楊生左肘。\par
路旁時賣故侯瓜，門前學種先生柳。\par
蒼茫古木連窮巷，寥落寒山對虛牖。\par
誓令疏勒出飛泉，不似潁川空使酒。\par
賀蘭山下陣如雲，羽檄交馳日夕聞。\par
節使三河募年少，詔書五道出將軍。\par
試拂鐵衣如雪色，聊持寶劍動星文。\par
願得燕弓射大將，恥令越甲鳴吾君。\par
莫嫌舊日雲中守，猶堪一戰取功勳。\par
桃源行·王維\par
漁舟逐水愛山春，兩岸桃花夾古津。\par
坐看紅樹不知遠，行盡青溪不見人。\par
山口潛行始隈隩，山開曠望旋平陸。\par
遙看一處攢雲樹，近入千家散花竹。\par
樵客初傳漢姓名，居人未改秦衣服。\par
居人共住武陵源，還從物外起田園。\par
月明松下房櫳靜，日出雲中雞犬喧。\par
驚聞俗客爭來集，競引還家問都邑。\par
平明閭巷掃花開，薄暮漁樵乘水入。\par
初因避地去人間，及至成仙遂不還。\par
峽裏誰知有人事，世中遙望空雲山。\par
不疑靈境難聞見，塵心未盡思鄉縣。\par
出洞無論隔山水，辭家終擬長遊衍。\par
自謂經過舊不迷，安知峯壑今來變。\par
當時只記入山深，青溪幾度到雲林。\par
春來遍是桃花水，不辨仙源何處尋。\par
清平調·雲想衣裳花想容·李白\par
雲想衣裳花想容，春風拂檻露華濃。\par
若非羣玉山頭見，會向瑤臺月下逢。\par
清平調·一枝紅豔露凝香·李白\par
一枝紅豔露凝香，雲雨巫山枉斷腸。\par
借問漢宮誰得似，可憐飛燕倚新妝。\par
清平調·名花傾國兩相歡·李白\par
名花傾國兩相歡，常得君王帶笑看。\par
解釋春風無限恨，沉香亭北倚欄杆。\par
行路難·金樽清酒鬥十千·李白\par
金樽清酒斗十千，玉盤珍羞直萬錢。(羞 通：饈；直 通 值)\par
停杯投箸不能食，拔劍四顧心茫然。\par
欲渡黃河冰塞川，將登太行雪滿山。(雪滿山 一作：雪暗天)\par
閒來垂釣碧溪上，忽復乘舟夢日邊。\par
行路難！行路難！多歧路，今安在？\par
長風破浪會有時，直掛雲帆濟滄海。\par
行路難·大道如青天·李白\par
大道如青天，我獨不得出。\par
羞逐長安社中兒，赤雞白雉賭梨慄。\par
彈劍作歌奏苦聲，曳裾王門不稱情。\par
淮陰市井笑韓信，漢朝公卿忌賈生。\par
君不見昔時燕家重郭隗，擁篲折節無嫌猜。\par
劇辛樂毅感恩分，輸肝剖膽效英才。\par
昭王白骨縈蔓草，誰人更掃黃金臺？\par
行路難，歸去來！\par
行路難·有耳莫洗潁川水·李白\par
有耳莫洗潁川水，有口莫食首陽蕨。\par
含光混世貴無名，何用孤高比雲月？\par
吾觀自古賢達人，功成不退皆殞身。\par
子胥既棄吳江上，屈原終投湘水濱。\par
陸機雄才豈自保？李斯稅駕苦不早。\par
華亭鶴唳詎可聞？上蔡蒼鷹何足道？\par
君不見吳中張翰稱達生，秋風忽憶江東行。\par
且樂生前一杯酒，何須身後千載名？\par
將進酒·李白\par
君不見，黃河之水天上來，奔流到海不復回。\par
君不見，高堂明鏡悲白髮，朝如青絲暮成雪。\par
人生得意須盡歡，莫使金樽空對月。\par
天生我材必有用，千金散盡還復來。\par
烹羊宰牛且爲樂，會須一飲三百杯。\par
岑夫子，丹丘生，將進酒，杯莫停。\par
與君歌一曲，請君爲我傾耳聽。(傾耳聽 一作：側耳聽)\par
鐘鼓饌玉不足貴，但願長醉不復醒。(不足貴 一作：何足貴；不復醒 一作：不願醒/不用醒)\par
古來聖賢皆寂寞，惟有飲者留其名。(古來 一作：自古；惟 通：唯)\par
陳王昔時宴平樂，斗酒十千恣歡謔。\par
主人何爲言少錢，徑須沽取對君酌。\par
五花馬，千金裘，呼兒將出換美酒，與爾同銷萬古愁。\par
玉階怨·李白\par
玉階生白露，夜久侵羅襪。\par
卻下水晶簾，玲瓏望秋月。(水晶 一作 水精)\par
長相思·其一·李白\par
長相思，在長安。\par
絡緯秋啼金井闌，微霜悽悽簟色寒。\par
孤燈不明思欲絕，卷帷望月空長嘆。\par
美人如花隔雲端！\par
上有青冥之長天，下有淥水之波瀾。\par
天長路遠魂飛苦，夢魂不到關山難。\par
長相思，摧心肝！\par
長相思·其二·李白\par
日色慾盡花含煙，月明欲素愁不眠。(欲素 一作：如素)\par
趙瑟初停鳳凰柱，蜀琴欲奏鴛鴦弦。\par
此曲有意無人傳，願隨春風寄燕然。\par
憶君迢迢隔青天，昔日橫波目，今作流淚泉。\par
不信妾斷腸，歸來看取明鏡前。(斷腸 一作：腸斷)\par
長幹行·其一·李白\par
妾發初覆額，折花門前劇。\par
郎騎竹馬來，繞牀弄青梅。\par
同居長幹裏，兩小無嫌猜，\par
十四爲君婦，羞顏未嘗開。\par
低頭向暗壁，千喚不一回。\par
十五始展眉，願同塵與灰。\par
常存抱柱信，豈上望夫臺。\par
十六君遠行，瞿塘灩澦堆。\par
五月不可觸，猿聲天上哀。(猿\par
一作：鳴)\par
門前遲行跡，一一生綠苔。\par
苔深不能掃，落葉秋風早。\par
八月蝴蝶來，雙飛西園草。\par
感此傷妾心，坐愁紅顏老。\par
早晚下三巴，預將書報家。\par
相迎不道遠，直至長風沙。\par
蜀道難·李白\par
噫籲嚱，危乎高哉！\par
蜀道之難，難於上青天！\par
蠶叢及魚鳧，開國何茫然！\par
爾來四萬八千歲，不與秦塞通人煙。\par
西當太白有鳥道，可以橫絕峨眉巔。\par
地崩山摧壯士死，然後天梯石棧相鉤連。\par
上有六龍回日之高標，下有衝波逆折之回川。\par
黃鶴之飛尚不得過，猿猱欲度愁攀援。(攀援 一作：攀緣)\par
青泥何盤盤，百步九折縈巖巒。\par
捫參歷井仰脅息，以手撫膺坐長嘆。\par
問君西遊何時還？畏途巉巖不可攀。\par
但見悲鳥號古木，雄飛雌從繞林間。\par
又聞子規啼夜月，愁空山。\par
蜀道之難,難於上青天，使人聽此凋朱顏！\par
連峯去天不盈尺，枯鬆倒掛倚絕壁。\par
飛湍瀑流爭喧豗，砯崖轉石萬壑雷。\par
其險也如此，嗟爾遠道之人胡爲乎來哉！(也如此 一作：也若此)\par
劍閣崢嶸而崔嵬，一夫當關，萬夫莫開。\par
所守或匪親，化爲狼與豺。\par
朝避猛虎，夕避長蛇；磨牙吮血，殺人如麻。\par
錦城雖雲樂，不如早還家。\par
蜀道之難,難於上青天，側身西望長諮嗟！\par
子夜吳歌·春歌·李白\par
秦地羅敷女，採桑綠水邊。\par
素手青條上，紅妝白日鮮。\par
蠶飢妾欲去，五馬莫留連。\par
子夜吳歌·夏歌·李白\par
鏡湖三百里，菡萏發荷花。\par
五月西施採，人看隘若耶。\par
回舟不待月，歸去越王家。\par
子夜吳歌·秋歌·李白\par
長安一片月，萬戶搗衣聲。\par
秋風吹不盡，總是玉關情。\par
何日平胡虜，良人罷遠征。\par
子夜吳歌·冬歌·李白\par
明朝驛使發，一夜絮徵袍。\par
素手抽針冷，那堪把剪刀。\par
裁縫寄遠道，幾日到臨洮。\par
關山月·李白\par
明月出天山，蒼茫雲海間。\par
長風幾萬裏，吹度玉門關。\par
漢下白登道，胡窺青海灣。\par
由來征戰地，不見有人還。\par
戍客望邊邑，思歸多苦顏。(望邊邑 一作：望邊色)\par
高樓當此夜，嘆息未應閒。\par
江南曲·李益\par
嫁得瞿塘賈，朝朝誤妾期。\par
早知潮有信，嫁與弄潮兒。\par
古從軍行·李頎\par
白日登山望烽火，黃昏飲馬傍交河。\par
行人刁斗風沙暗，公主琵琶幽怨多。\par
野雲萬里無城郭，雨雪紛紛連大漠。\par
胡雁哀鳴夜夜飛，胡兒眼淚雙雙落。\par
聞道玉門猶被遮，應將性命逐輕車。\par
年年戰骨埋荒外，空見蒲桃入漢家。\par
哀王孫·杜甫\par
長安城頭頭白烏，夜飛延秋門上呼。\par
又向人家啄大屋，屋底達官走避胡。\par
金鞭斷折九馬死，骨肉不得同馳驅。\par
腰下寶玦青珊瑚，可憐王孫泣路隅。\par
問之不肯道姓名，但道困苦乞爲奴。\par
已經百日竄荊棘，身上無有完肌膚。\par
高帝子孫盡隆準，龍種自與常人殊。\par
豺狼在邑龍在野，王孫善保千金軀。\par
不敢長語臨交衢，且爲王孫立斯須。\par
昨夜東風吹血腥，東來橐駝滿舊都。\par
朔方健兒好身手，昔何勇銳今何愚。\par
竊聞天子已傳位，聖德北服南單于。\par
花門剺面請雪恥，慎勿出口他人狙。\par
哀哉王孫慎勿疏，五陵佳氣無時無。\par
兵車行·杜甫\par
車轔轔，馬蕭蕭，行人弓箭各在腰。\par
耶孃妻子走相送，塵埃不見咸陽橋。(\par
娘 一作：“爺”)\par
牽衣頓足攔道哭，哭聲直上幹雲霄。\par
道旁過者問行人，行人但云點行頻。\par
或從十五北防河，便至四十西營田。\par
去時里正與裹頭，歸來頭白還戍邊。\par
邊庭流血成海水，武皇開邊意未已。\par
君不聞，漢家山東二百州，千村萬落生荊杞。\par
縱有健婦把鋤犁，禾生隴畝無東西。\par
況復秦兵耐苦戰，被驅不異犬與雞。\par
長者雖有問，役夫敢申恨？\par
且如今年冬，未休關西卒。\par
縣官急索租，租稅從何出？\par
信知生男惡，反是生女好。\par
生女猶得嫁比鄰，生男埋沒隨百草。\par
君不見，青海頭，古來白骨無人收。\par
新鬼煩冤舊鬼哭，天陰雨溼聲啾啾！\par
麗人行·杜甫\par
三月三日天氣新，長安水邊多麗人。\par
態濃意遠淑且真，肌理細膩骨肉勻。\par
繡羅衣裳照暮春，蹙金孔雀銀麒麟。\par
頭上何所有？翠微盍葉垂鬢脣。\par
背後何所見？珠壓腰衱穩稱身。\par
就中雲幕椒房親，賜名大國虢與秦。\par
紫駝之峯出翠釜，水精之盤行素鱗。\par
犀箸厭飫久未下，鸞刀縷切空紛綸。\par
黃門飛鞚不動塵，御廚絡繹送八珍。\par
簫鼓哀吟感鬼神，賓從雜遝實要津。\par
後來鞍馬何逡巡，當軒下馬入錦茵。\par
楊花雪落覆白蘋，青鳥飛去銜紅巾。\par
炙手可熱勢絕倫，慎莫近前丞相嗔！\par
哀江頭·杜甫\par
少陵野老吞聲哭，春日潛行曲江曲。\par
江頭宮殿鎖千門，細柳新蒲爲誰綠？\par
憶昔霓旌下南苑，苑中萬物生顏色。\par
昭陽殿裏第一人，同輦隨君侍君側。\par
輦前才人帶弓箭，白馬嚼齧黃金勒。\par
翻身向天仰射雲，一笑正墜雙飛翼。\par
明眸皓齒今何在？血污遊魂歸不得。\par
清渭東流劍閣深，去住彼此無消息。\par
人生有情淚沾臆，江水江花豈終極！\par
黃昏胡騎塵滿城，欲往城南望城北。\par
金縷衣·佚名\par
勸君莫惜金縷衣，勸君惜取少年時。(惜取 一作：須取)\par
花開堪折直須折，莫待無花空折枝。(花開 一作：有花)\par
獨不見·沈佺期\par
盧家少婦鬱金堂，海燕雙棲玳瑁梁。\par
九月寒砧催木葉，十年征戍憶遼陽。\par
白狼河北音書斷，丹鳳城南秋夜長。\par
誰謂含愁獨不見，更教明月照流黃。\par
烈女操·孟郊\par
梧桐相待老，鴛鴦會雙死。\par
貞女貴殉夫，捨生亦如此。\par
波瀾誓不起，妾心古井水。(古井水 一作：井中水)\par
遊子吟·孟郊\par
慈母手中線，遊子身上衣。\par
臨行密密縫，意恐遲遲歸。\par
誰言寸草心，報得三春暉。\par
燕歌行·高適\par
漢家煙塵在東北，漢將辭家破殘賊。\par
男兒本自重橫行，天子非常賜顏色。\par
摐金伐鼓下榆關，旌旆逶迤碣石間。\par
校尉羽書飛瀚海，單于獵火照狼山。\par
山川蕭條極邊土，胡騎憑陵雜風雨。\par
戰士軍前半死生，美人帳下猶歌舞。\par
大漠窮秋塞草腓，孤城落日鬥兵稀。\par
身當恩遇恆輕敵，力盡關山未解圍。\par
鐵衣遠戍辛勤久，玉箸應啼別離後。\par
少婦城南欲斷腸，徵人薊北空回首。\par
邊庭飄颻那可度，絕域蒼茫更何有。\par
殺氣三時作陣雲，寒聲一夜傳刁斗。\par
相看白刃血紛紛，死節從來豈顧勳。\par
君不見沙場征戰苦，至今猶憶李將軍。\par
長幹行·君家何處住·崔顥\par
君家何處住，妾住在橫塘。\par
停船暫借問，或恐是同鄉。\par
長幹行·家臨九江水·崔顥\par
家臨九江水，來去九江側。\par
同是長幹人，生小不相識。\par
塞下曲·鷲翎金僕姑·盧綸\par
鷲翎金僕姑，燕尾繡蝥弧。\par
獨立揚新令，千營共一呼。\par
塞下曲·林暗草驚風·盧綸\par
林暗草驚風，將軍夜引弓。\par
平明尋白羽，沒在石棱中。\par
塞下曲·月黑雁飛高·盧綸\par
月黑雁飛高，單于夜遁逃。\par
欲將輕騎逐，大雪滿弓刀。\par
塞下曲·野幕敞瓊筵·盧綸\par
野幕敞瓊筵，羌戎賀勞旋。\par
醉和金甲舞，雷鼓動山川。\par

\end{document}